\documentclass[12pt,b5paper,twoside]{article}

\usepackage[T1]{fontenc}
\usepackage[utf8]{inputenc}
\usepackage[english]{babel}
\usepackage{textcomp}
\usepackage{newpxtext}
\usepackage{newpxmath}
\usepackage{mathtools}
\usepackage{enumitem}
\usepackage{microtype}
\usepackage[b5paper,
  left=0.7cm,
  right=0.7cm,
  top=0.91cm,
  bottom=0.91cm,
  includehead,
  headheight=16pt,
  headsep=2pt,
  footskip=6pt
]{geometry}
\usepackage{fancyhdr}
\usepackage{xcolor}
\definecolor{InternalLink}{HTML}{245B91}
\definecolor{CitationLink}{HTML}{7A3E9D}
\definecolor{ExternalLink}{HTML}{007C83}
\usepackage[
  unicode=true,
  colorlinks=true,
  linkcolor=InternalLink,
  citecolor=CitationLink,
  urlcolor=ExternalLink,
  pdfborder={0 0 0}
]{hyperref}

\hypersetup{
  pdftitle={Über Funktionen von Funktionaloperatoren / On Functions of Functional Operators},
  pdfauthor={J. von Neumann},
  pdfsubject={Faithful English translation and LaTeX edition}
}

\makeatletter
\renewcommand\normalsize{%
  \@setfontsize\normalsize{13pt}{15.6pt}%
  \abovedisplayskip 10pt plus 2pt minus 4pt%
  \belowdisplayskip \abovedisplayskip%
  \abovedisplayshortskip 6pt plus 2pt%
  \belowdisplayshortskip 6pt plus 2pt minus 2pt}
\renewcommand\small{\@setfontsize\small{11.5pt}{13.8pt}}
\renewcommand\footnotesize{\@setfontsize\footnotesize{10.5pt}{12.6pt}}
\renewcommand\scriptsize{\@setfontsize\scriptsize{9.5pt}{11.4pt}}
\makeatother

\linespread{1.1}
\allowdisplaybreaks[4]
\setlength{\parindent}{1.35em}
\setlength{\parskip}{0pt}
\setlength{\fboxsep}{2pt}
\renewcommand{\footnoterule}{}
\emergencystretch=1.5em
\clubpenalty=8000
\widowpenalty=8000
\displaywidowpenalty=8000

\pagestyle{fancy}
\fancyhf{}
\fancyhead[LE]{\thepage}
\fancyhead[RO]{\thepage}
\renewcommand{\headrulewidth}{0pt}
\renewcommand{\footrulewidth}{0pt}

\newcommand{\noteref}[1]{%
  \hyperlink{note:#1}{\textsuperscript{#1)}}}
\newcommand{\originalnote}[1]{%
  \par\smallskip\noindent\hypertarget{note:#1}{\textsuperscript{\textbf{#1)}}}\,}
\newcommand{\tnref}[1]{%
  \hyperlink{tn:#1}{\textsuperscript{\textsc{TN-#1}}}}
\newcommand{\translatornote}[1]{%
  \par\smallskip\noindent\hypertarget{tn:#1}{\textbf{TN-#1.}}\,}
\newcommand{\TheoremHead}[1]{%
  \par\smallskip\noindent\textbf{Theorem #1.}\,}
\newcommand{\DefinitionHead}{%
  \par\smallskip\noindent\textbf{Definition.}\,}
\newcommand{\ProofHead}{%
  \par\smallskip\noindent\textit{Proof.}\,}

\begin{document}
\normalsize
\setcounter{page}{1}
\thispagestyle{empty}

\begin{center}
  {\fontsize{17pt}{20.4pt}\selectfont\bfseries
    ÜBER FUNKTIONEN VON FUNKTIONALOPERATOREN%
    \begingroup\renewcommand{\thefootnote}{*}\footnotemark\endgroup\par}
  \vspace{0.35em}
  {\fontsize{17pt}{20.4pt}\selectfont\bfseries
    ON FUNCTIONS OF FUNCTIONAL OPERATORS\par}
  \vspace{0.7em}
  {\large J. von Neumann, Berlin.\par}
\end{center}
\begingroup
  \renewcommand{\thefootnote}{*}%
  \footnotetext{Received October 20, 1930.}%
\endgroup

% Source PDF page 1 / journal p. 191
\section*{Introduction}
\label{sec:intro}

\subsection*{1.}
\label{sec:intro.1}
The subject of the present paper is formed by certain theorems concerning the functional dependence of commuting (bounded) functional operators. They are connected with questions already touched upon in an earlier paper by the author\noteref{1}, but not discussed there in full (for reasons of space, since that paper was also devoted to other subjects). Here an exhaustive investigation and resolution will be undertaken, together with a complete development, suggested by it, of the concept of a function for functional operators. What is at issue is the proof of the following theorem: two (or even arbitrarily many) commuting\noteref{2} Hermitian operators can always be represented as functions of a single one. In A, as cited in \noteref{1}, it was shown only that, in a certain sense, they are limits of polynomials in that same operator, since the general concept of a function of an operator was not developed there. The theorem just stated will, moreover, be proved here in a somewhat strengthened and more precise form, to which we shall return below.

\subsection*{2.}
\label{sec:intro.2}
Among the concepts essential to the following considerations, those of Hilbert space and of its linear operators occupy first place. It is probably unnecessary to develop them in detail; but since in investigating them we shall employ a geometric terminology specially adapted to our purposes, let it be said that we follow the relevant exposition and definitions in E and A\noteref{3}. As there, we denote the (abstract) Hilbert space by $\mathfrak H$, and the ring of its bounded operators by $\mathbf B$. We shall also repeatedly speak of rings in $\mathbf B$: these are those subsets of $\mathbf B$ which, together with any elements $A,B$, contain $aA$, $A^*$\noteref{4}, $A+B$, and $AB$, and which in addition possess a certain closure property\noteref{5}. Every subset $\mathbf M$ of $\mathbf B$ is contained in certain rings, among them a smallest one: this is the ring $\mathbf R(\mathbf M)$ generated by it\noteref{6}. In particular, if $\mathbf M$ consists of a single element $A$, we can generate from $A$ the ring $\mathbf R(A)$. This consists

% Source PDF page 2 / journal p. 192
of all polynomials $p(A)$ in $A$ (where $p(x)$ ranges over the polynomials satisfying $p(0)=0$), together with the limits of all so-called strongly doubly convergent sequences of such polynomials\noteref{7}--at least when $A$ is Hermitian, or at any rate normal\noteref{8}.

A set $\mathbf M$ (in $\mathbf B$) is Abelian if all its elements commute with one another and with one another's adjoints; for sets which contain $A^*$ whenever they contain $A$, for example rings, pairwise commutativity of all elements therefore suffices\noteref{9}. It was also shown in A, at the place cited in \noteref{9}, that $\mathbf M$ is Abelian if and only if the ring $\mathbf R(\mathbf M)$ is Abelian, and that a ring is Abelian if and only if it consists entirely of normal operators.

\subsection*{3.}
\label{sec:intro.3}
It was proved in A\noteref{10} that, as $A$ ranges over all normal operators, $\mathbf R(A)$ ranges over all Abelian rings; indeed, one may restrict $A$ to the Hermitian operators (cf. \noteref{8}). That is, the elements of the ring are, in the sense sketched above, limits of polynomials in $A$. We wish, however, to represent them directly as functions of $A$. To this end we shall introduce in this paper the general concept of a function of a Hermitian (or normal) operator (and delimit precisely the class of functions admissible for this purpose), and shall then prove that $\mathbf R(A)$ is the set of all $f(A)$, where $f(x)$ ranges over all admissible functions satisfying $f(0)=0$; in fact, $f(x)$ may be restricted to functions of the first Baire class.

Since every Abelian set is part of an Abelian ring (for example, the ring it generates), its elements are functions of a single $A$. Every set of commuting Hermitian operators is Abelian by definition (because $B^*=B$ for each of its elements), and the same assertion therefore holds for it as well. We shall analyze the various possible formulations and strengthenings of this theorem more closely below.

\subsection*{4.}
\label{sec:intro.4}
As is clear from what has been said, our principal task is to introduce the general concept of a function of an operator. There are two ways to do this: one may proceed by successive passages to the limit for functions of all Baire classes, or one may give a definition by means of Lebesgue--Stieltjes integrals, at one stroke, for the (more general) full class of functions for which this is possible at all. These two methods are, moreover, most closely related to two routes for introducing the general concept of integration, those of Young and Lebesgue, respectively\noteref{11}. We shall use the second, which permits more general results, although the first would also suffice to establish the theorem stated in \hyperref[sec:intro.3]{paragraph 3}.

Before approaching our proper subject, however, we must prove an auxiliary theorem about Lebesgue integrals which, together with some consequences, may perhaps also be not entirely without

% Source PDF page 3 / journal p. 193
interest in its own right and which, so far as the author is aware, has not previously appeared in the literature.

The paper is organized as follows. In \hyperref[sec:I]{Chapter I} we prove the auxiliary theorems about Lebesgue integrals just mentioned and then introduce the Lebesgue--Stieltjes integral, together with those of its properties that are essential for our purposes. (A short, fully rigorous digression suffices for this, and, going somewhat beyond the customary manner of definition, we shall also need it for discontinuous differentials; the exposition is therefore included although it offers nothing essentially new.) In \hyperref[sec:II]{Chapter II} we define functions of operators and establish their principal properties; in Chapter III the theorem mentioned in \hyperref[sec:intro.3]{paragraph 3}, together with several consequences, is proved. In Chapter IV we generalize the concept of a function--which up to that point is available only for a single argument, required to be Hermitian--to finitely or countably infinitely many arguments, which may be commuting normal operators.

\section*{I. Auxiliary Theorems on Lebesgue and Lebesgue--Stieltjes Integrals}
\label{sec:I}

\subsection*{1.}
\label{sec:I.1}
Let $\varphi(a)$ be a function defined on an interval $0\leq a\leq\alpha$ (where $\alpha$ is any number $\geq0$), monotonically nondecreasing, never negative, and everywhere right-continuous; for short, we shall call such a function an \emph{$M$-function}.

Let $f(x)$ be a measurable, never negative, bounded function defined on a measurable set $\mathfrak M$ of finite measure. We define the measure function $\varphi(a)$ of $f(x)$ as follows:
\[
  \varphi(a)=\operatorname{meas}\{x\in\mathfrak M:f(x)\leq a\}.
\]
Evidently $\varphi(a)$ is an $M$-function; for $\alpha$ we take the supremum of the range of $f(x)$.

Since every $M$-function $\varphi(a)$ itself possesses the properties of the function $f(x)$ just listed, we may in turn form its measure function, the $M$-function $\psi(b)$. The following relation between $\varphi(a)$ and $\psi(b)$ is evident (and may also be used as a definition):
\[
  \psi(b)=\sup\{a:0\leq a<\alpha,\ \varphi(a)\leq b\}.
\]
The domain of $\psi(b)$ is $0\leq b\leq\varphi(\alpha)$. From the preceding relation one readily deduces that the measure function of $\psi(b)$ is again $\varphi(a)$; that is, the relation between $\varphi(a)$ and $\psi(b)$ is reciprocal. Furthermore, $\varphi(\psi(b))=b$ whenever $b$ is not at the left endpoint or in the interior of an interval of constancy of $\psi$ (that is, at the lower endpoint or in the interior of an interval produced, at a discontinuity of $\varphi$, by its ``jump''); and likewise $\psi(\varphi(a))=a$, with the corresponding exception for $a$. Thus $\varphi(a)$ and $\psi(b)$ are inverse to one another insofar as this is possible at all.

% Source PDF page 4 / journal p. 194
\subsection*{2.}
\label{sec:I.2}
Let $f(x)$ be as at the beginning of \hyperref[sec:I.1]{subsection 1}, let $\varphi(a)$ be its measure function, and let $\psi(b)$ be the measure function of $\varphi(a)$. If $f(x)$ is an $M$-function, then, as we have seen, $f=\psi$ (and only then, since $\psi(b)$ is an $M$-function); but for arbitrary $f(x)$ as well, there is a very close relation between it and $\psi(b)$. Namely, we shall prove the following.

\phantomsection\label{thm:1}
\noindent\textbf{Theorem 1.} \emph{Let $g(u)$ be an arbitrary function\noteref{12}. Then}
\begin{equation}
  \int_{\mathfrak M} g(f(x))\,dx
  =\int_{0}^{\operatorname{meas}\mathfrak M} g(\psi(b))\,db,
  \label{eq:thm1-equiintegrability}
\end{equation}
\emph{where $0\leq b\leq\operatorname{meas}\mathfrak M$ is the domain of $\psi(b)$, since $\operatorname{meas}\mathfrak M$ is the maximum of $\varphi(a)$. The equality is to be understood in the sense that the left-hand side is always meaningful whenever the right-hand side is meaningful (this amounts to Lebesgue integrability of the respective integrands), and the two are then equal.}

\textit{Proof.} It is enough to show that if all the sets of $b$ for which $g(\psi(b))\leq m$ are measurable, then all the sets of $x$ for which $g(f(x))\leq m$ are measurable as well, and that the respective sets have equal measures. This is certainly established if the following holds for every set $\mathfrak P$: whenever the set of $b$ for which $\psi(b)$ belongs to $\mathfrak P$ is measurable, so too is the set of $x$ for which $f(x)$ belongs to $\mathfrak P$, and the two have the same measure. (One then takes $\mathfrak P$ to be the set of $u$ for which $g(u)\leq m$.) Call the set of $b$ in question $\mathfrak P_b$, and the set of $x$ in question $\mathfrak P_x$.

We now use $\mathfrak P_b$ as the point of departure. It necessarily has the following property: if $I$ is an interval of constancy of $\psi(b)$\noteref{13}, then either $I$ is contained in $\mathfrak P_b$ or it is disjoint from $\mathfrak P_b$. The number of the mutually disjoint intervals of constancy of $\psi(b)$ is finite or countable; denote them by $I_1,I_2,\ldots$. We now regard $\mathfrak P_b$, independently of $\mathfrak P$, as an arbitrary set, subject only to the requirement that every $I_n$ be either contained in $\mathfrak P_b$ or disjoint from it (and that $\mathfrak P_b$ be contained in the interval $0\leq b\leq\operatorname{meas}\mathfrak M$). The set $\mathfrak P_x$ consists of all $x\in\mathfrak M$ for which $f(x)=\psi(b)$ for a suitable $b\in\mathfrak P_b$. Thus it remains to show that if $\mathfrak P_b$ is measurable, then $\mathfrak P_x$ is measurable too and the two sets have the same measure.

We shall show that it suffices to prove this for those $\mathfrak P_b$ that are intervals $0\leq b<\beta$ or $0\leq b\leq\beta$ (and that stand in the relation just specified to the $I_n$). From these it passes first to their differences, hence to all intervals (with arbitrary assignment of endpoints, still subject to the stated relation to the $I_n$), and therefore in particular to open intervals and to the $I_n$. From these it passes in turn to unions of finitely or countably many such sets: first, therefore, to all unions of arbitrary $I_n$, and second to all open sets (that is, sets relatively open in $0\leq b\leq\operatorname{meas}\mathfrak M$, naturally still subject to the stated relation to the $I_n$). In particular, it holds for the full interval $0\leq b\leq\operatorname{meas}\mathfrak M$.

% Source PDF page 5 / journal p. 195
It now suffices to prove, for every one of our sets $\mathfrak P_b$, that
\[
  \operatorname{meas}\mathfrak P_b
  \geq \operatorname{outer\,meas}\mathfrak P_x
\]
(without assuming that $\mathfrak P_x$ is measurable). Indeed, on applying this to the complement of $\mathfrak P_b$, and using what was just proved for the full interval $0\leq b\leq\operatorname{meas}\mathfrak M$, one obtains
\[
  \operatorname{meas}\mathfrak P_b
  \leq \operatorname{inner\,meas}\mathfrak P_x;
\]
hence
\[
  \operatorname{meas}\mathfrak P_b
  =\operatorname{outer\,meas}\mathfrak P_x
  =\operatorname{inner\,meas}\mathfrak P_x.
\]

Let $\mathfrak Q_b$ be the union of all $I_n$ disjoint from $\mathfrak P_b$. Then $\operatorname{meas}\mathfrak Q_b=\operatorname{meas}\mathfrak Q_x$, and it is consequently enough to prove
\[
 \operatorname{meas}(\mathfrak Q_b\cup\mathfrak P_b)
 \geq
 \operatorname{outer\,meas}(\mathfrak Q_x\cup\mathfrak P_x).
\]
The set $\mathfrak Q_b\cup\mathfrak P_b$ contains every $I_n$; if we denote it again by $\mathfrak P_b$, we find that it suffices to consider the preceding relation for those $\mathfrak P_b$ which contain all the $I_n$.

Now $\mathfrak P_b$ is contained in an open set $\mathfrak P'_b$ whose measure exceeds that of $\mathfrak P_b$ by at most $\varepsilon$ (and this can be done for every $\varepsilon>0$). It therefore suffices to consider open $\mathfrak P'_b$ containing all the $I_n$: if our relation holds for these, it also holds for $\mathfrak P_b$ itself. For open $\mathfrak P'_b$ (which do indeed stand in the previously required relation to every $I_n$\noteref{14}), however, we already know the relation, even with equality.

Thus it does in fact suffice to consider the intervals $\mathfrak P_b:0\leq b<\beta$ and $0\leq b\leq\beta$. Their relation to the $I_n$ now implies that each really is the $b$-set corresponding to some $\mathfrak P$ (that is, $b\in\mathfrak P_b$ when $\psi(b)\in\mathfrak P$); by monotonicity of $\psi(b)$, $\mathfrak P$ is an initial interval of $0\leq u<\infty$, just as $\mathfrak P_b$ is. Thus $\mathfrak P$ is a set $0\leq u<a$ or $0\leq u\leq a$. Since $\psi(b)$ and $f(x)$ are never negative, the condition $0\leq u$ may be omitted. Since our assertion (that $\mathfrak P_b$ and $\mathfrak P_x$ are measurable and have the same measure) passes immediately from an increasing sequence of $\mathfrak P$-sets to their union, it is enough to consider $u\leq a$. Hence $\mathfrak P_b$ is defined by $\psi(b)\leq a$, and $\mathfrak P_x$ by $f(x)\leq a$. Both sets are plainly measurable; and since $\psi(b)$ and $f(x)$ have the same measure function $\varphi(a)$, both have the same measure, namely $\varphi(a)$.

This completes the proof.

\subsection*{3.}
\label{sec:I.3}
We shall draw several consequences from the result formulated in \hyperref[thm:1]{Theorem 1}. For this purpose we introduce the following notion. If a function $f(x)$ (defined for every real $x$ and having complex values) is such that $f(\Phi(a))$ is measurable for every monotonically nondecreasing, right-continuous function $\Phi(a)$, we call $f$ \emph{$\Phi$-measurable}\noteref{15}. Evidently all polynomials are $\Phi$-measurable. If $f_1,f_2,\ldots$ is a sequence of $\Phi$-measurable functions converging to $f$ (that is, $f_n(x)\to f(x)$ for every $x$ as $n\to\infty$), then $f$ is also $\Phi$-measurable, by the analogous property of ordinary measurability. Thus all functions of all Baire classes are $\Phi$-measurable. It is also clear that, together with $f(x)$ and $g(x)$, the functions $f(x)\pm g(x)$, $af(x)$, and $f(x)g(x)$ are $\Phi$-measurable, again by the corresponding property of ordinary

% Source PDF page 6 / journal p. 196
measurability. The following fact, however, is by no means trivial.

\phantomsection\label{thm:2}
\noindent\textbf{Theorem 2.} \emph{If $g(u)$ is $\Phi$-measurable and $f(x)$ is measurable, respectively $\Phi$-measurable (but real-valued), then $g(f(x))$ is measurable, respectively $\Phi$-measurable.}

\textit{Proof.} It suffices to consider measurability, since the case of $\Phi$-measurability is reduced to it by replacing $f(x)$ by $f(\Phi(a))$. We may further suppose that $g(u)$ is bounded, since otherwise we may replace it, for example, by $\operatorname{Tghyp}(g(u))$. The range of $f(x)$ may be assumed to lie in the interval $0<u<2$, for otherwise we replace $f(x)$ by $1+\operatorname{Tghyp}(f(x))$ and $g(u)$ by
\[
 \overline g(u)=g\bigl(\operatorname{Arctghyp}(u-1)\bigr)
 \quad(0<u<2),
\]
taking, say, $\overline g(u)=0$ otherwise. Finally, the domain of $f(x)$ may also be taken to be $0<x<2$, since otherwise we consider $f(\operatorname{Arctghyp}(x-1))$.

Thus $g(u)$ is bounded, while $f(x)$ is nonnegative, bounded, and defined on a set of finite measure. Let $\Phi(a)$ be the measure function of the measure function of $f(x)$. Then $g(\Phi(a))$ is measurable and hence integrable. By \hyperref[thm:1]{Theorem 1}, $g(f(x))$ is therefore integrable, and hence measurable.

A second fact, important for later use, is the following.

\phantomsection\label{thm:3}
\noindent\textbf{Theorem 3.} \emph{Let $F(u)$ be a $\Phi$-measurable function, and let $f_1(x),f_2(x),\ldots$ be arbitrary measurable functions. Then there exists a function $G(u)$ of the third Baire class such that the set $\mathfrak P$ of all $u$ for which $F(u)\neq G(u)$ has the following property: the set of all $x$ for which at least one $f_n(x)$ belongs to $\mathfrak P$ is a Lebesgue null set\noteref{16}.}

\textit{Proof.} By replacing $f_n(x)$ with $f_n(\operatorname{Arctghyp}x)$ (the map $x\mapsto\operatorname{Arctghyp}x$ maps $-1<x<1$ onto the whole real line and carries Lebesgue null sets into Lebesgue null sets), we make its domain $-1<x<1$. By then replacing it with
\[
 f_n\!\left(\frac{2}{\beta-\alpha}
       \left(x-\frac{\alpha+\beta}{2}\right)\right)
\]
($\alpha<\beta$, with both numbers otherwise arbitrary), we make $\alpha<x<\beta$ its domain. We may therefore assume from the outset that $f_n(x)$ has the domain
\[
 \frac{1}{n+1}<x<\frac{1}{n}.
\]
All these functions $f_n(x)$, $n=1,2,\ldots$, can now be combined into a single measurable function $f(x)$ defined on $0<x<1$ (at $x=\tfrac12,\tfrac13,\ldots$, put, say, $f(x)=0$). Thus the problem is to choose $G(u)$ so that the set of $x$ for which $f(x)\in\mathfrak P$ is a Lebesgue null set. Let $\psi(b)$ be the measure function of the measure function of $f(x)$. By \hyperref[thm:1]{Theorem 1}, the set of $x$ for which $f(x)\in\mathfrak P$ is certainly a Lebesgue null set whenever the set of $b$ for which $\psi(b)\in\mathfrak P$ is one (take $g(u)=1$ for $u\in\mathfrak P$ and $g(u)=0$ otherwise). We may therefore replace $f(x)$ by $\psi(b)$.

% Source PDF page 7 / journal p. 197
That is, it suffices to consider the case in which only a single function $f_1(x)$ occurs and this function is an $M$-function (cf. \noteref{16}). Denote it by $\varphi(x)$ (in place of $\psi(b)$).

First, $F(u)$ is $\Phi$-measurable, so $F(\varphi(x))$ is measurable. By Vitali's theorem (cf. \noteref{16}), there is consequently a function $\overline G(x)$ of the second Baire class such that
\[
 F(\varphi(x))=\overline G(x)
\]
except on a Lebesgue null set. If $\overline G(x)$ had the form $G(\varphi(x))$, with $G(u)$ of the third Baire class, the proof would be complete. It is also complete if we can alter $\overline G(x)$ on a Lebesgue null set so as to obtain a function
\[
 \overline{\overline G}(x)=G(\varphi(x)).
\]
We shall now do this.

For $\overline{\overline G}(x)$ to have the form $G(\varphi(x))$ (for some $G(u)$) means simply that it is constant on every interval of constancy of $\varphi(x)$. The function $F(\varphi(x))$ has this property and agrees with $\overline G(x)$ except on a Lebesgue null set. Thus, by making $\overline G(x)$ constant on the intervals of constancy of $\varphi(x)$, we change it only on a Lebesgue null set. This gives a function $\overline{\overline G}(x)=G(\varphi(x))$; it remains only to show that $G(u)$ is of the third Baire class. In fact, if $\psi(u)$ denotes the measure function of $\varphi(x)$, we may put
\[
 G(u)=\overline{\overline G}(\psi(u)),
\]
because $\psi(\varphi(x))=x$ when $x$ does not belong to an interval of constancy of $\varphi(x)$, while if it does belong to one, then $\psi(\varphi(x))$ and $x$ lie in that same interval, on which $\overline{\overline G}(x)$ is constant. Hence $\overline{\overline G}(x)=G(\varphi(x))$.

First, $\overline{\overline G}(x)$ is of the second Baire class: it is obtained from a function $\overline G(x)$ of that class by altering it, namely by making it constant, on finitely or countably many intervals (the intervals of constancy of $\varphi(x)$). This assertion for the second Baire class is plainly valid in general if, for finitely many intervals, it is valid for the first Baire class\noteref{17}. A function of the first Baire class is the limit of an everywhere convergent sequence of continuous functions; when it is altered in the stated manner, it is the limit of the sequence of continuous functions altered in the same way. After this alteration, however, each of these functions is still continuous everywhere except at finitely many points (the endpoints of the intervals on which the alteration was made). It is easy to see that limits of such functions can themselves be represented as limits of continuous functions\noteref{18}. Finally, since $\overline{\overline G}(x)$ belongs to the second Baire class, we must show that $G(u)=\overline{\overline G}(\psi(u))$ belongs to the third. Now $\overline{\overline G}(x)$ is obtained by two successive passages to the limit from certain continuous functions $f_{m,n}(x)$\noteref{19}; hence $G(u)=\overline{\overline G}(\psi(u))$ is obtained in the same way from the functions $f_{m,n}(\psi(u))$, which, like $\psi(u)$, are right-continuous. Since every right-continuous function belongs to the first Baire class\noteref{20}, $G(u)$ belongs to the third.

% Source PDF page 8 / journal p. 198
This concludes our considerations concerning Lebesgue measure and the Lebesgue integral.

\subsection*{4.}
\label{sec:I.4}
We now define the Lebesgue--Stieltjes integral $\int f(x)\,d\varphi(x)$, initially for the case in which $\varphi(x)$ is an $M$-function. The domain of $\varphi(x)$ is then an interval $0\leq x\leq\alpha$, and these numbers shall also be the limits of integration; for the moment we keep them fixed. Let $\psi(y)$ be the measure function of $\varphi(x)$. Following Lebesgue, we define
\begin{equation}
 \int f(x)\,d\varphi(x)=\int f(\psi(y))\,dy,
 \label{eq:LS-definition-finite}
\end{equation}
where the right-hand side is an ordinary Lebesgue integral. It is certainly meaningful when $f(x)$ is bounded and $\Phi$-measurable; this will therefore always be assumed of $f(x)$ in what follows.

The integral so defined has the following properties:
\begin{description}
 \item[\textup{a)}]\phantomsection\label{prop:I.4.a}
 \[
  \int af(x)\,d\varphi(x)=a\int f(x)\,d\varphi(x)
 \]
 for every complex number $a$.

 \item[\textup{b)}]\phantomsection\label{prop:I.4.b}
 \[
  \int \bigl(f(x)+g(x)\bigr)\,d\varphi(x)
  =\int f(x)\,d\varphi(x)+\int g(x)\,d\varphi(x).
 \]

 \item[\textup{c)}]\phantomsection\label{prop:I.4.c}
 \[
  \int f(x)\,d\bigl(a\varphi(x)\bigr)
  =a\int f(x)\,d\varphi(x)\qquad(a>0).
 \]

 \item[\textup{d)}]\phantomsection\label{prop:I.4.d}
 \[
  \int f(x)\,d\bigl(\varphi(x)+\chi(x)\bigr)
  =\int f(x)\,d\varphi(x)+\int f(x)\,d\chi(x).
 \]

 \item[\textup{e)}]\phantomsection\label{prop:I.4.e}
 Suppose that $f(x)\geq0$ on a set $\mathfrak M$, and that the set of $y$ for which $\psi(y)\notin\mathfrak M$ is a Lebesgue null set\noteref{21}. Then
 \[
  \int f(x)\,d\varphi(x)\geq0.
 \]

 \item[\textup{f)}]\phantomsection\label{prop:I.4.f}
 If $f(x)=g(x)$ on a set $\mathfrak M$ as in \hyperref[prop:I.4.e]{e)}, then
 \[
  \int f(x)\,d\varphi(x)=\int g(x)\,d\varphi(x).
 \]

 \item[\textup{g)}]\phantomsection\label{prop:I.4.g}
 If $f_n(x)\to f(x)$ as $n\to\infty$ on a set $\mathfrak M$ as in \hyperref[prop:I.4.e]{e)}, and if the functions $f_n(x)$ are uniformly bounded, then
 \[
  \int f_n(x)\,d\varphi(x)\longrightarrow\int f(x)\,d\varphi(x).
 \]
\end{description}

Properties \hyperref[prop:I.4.a]{a)} and \hyperref[prop:I.4.b]{b)} follow at once from the corresponding properties of the Lebesgue integral, while \hyperref[prop:I.4.e]{e)}--\hyperref[prop:I.4.g]{g)} likewise pass immediately, by the definition, into familiar properties of that integral. Only \hyperref[prop:I.4.c]{c)} and \hyperref[prop:I.4.d]{d)} remain to be verified. By \hyperref[prop:I.4.f]{f)} and \hyperref[thm:3]{Theorem 3}, however, it is enough to prove them for $f(x)$ and $g(x)$ of the third Baire class, and by \hyperref[prop:I.4.g]{g)} even for those of the zeroth class, that is, for continuous $f(x)$ and $g(x)$. In that case we plainly have to do with ordinary

% Source PDF page 9 / journal p. 199
Riemann--Stieltjes integrals (for $\int f(\psi(y))\,dy$ is a Riemann integral, since $f(\psi(y))$ is right-continuous), and \hyperref[prop:I.4.c]{c)} and \hyperref[prop:I.4.d]{d)} are well known for these.

We also observe the following. Suppose that a continuous, monotonically increasing, and hence one-to-one map
\[
 \xi=\omega(x),\qquad x=\omega^{-1}(\xi),
\]
maps $0\leq x\leq\alpha$ onto $0\leq\xi\leq\beta$. Then
\begin{equation}
 \int f(\xi)\,d\varphi(\xi)
 =\int f(\omega(x))\,d\varphi(\omega(x)).
 \label{eq:LS-change-finite}
\end{equation}
(That $f(\omega(x))$ is $\Phi$-measurable follows immediately from the $\Phi$-measurability of $f(\xi)$.) Indeed, if $\psi(y)$ is the measure function of $\varphi(x)$, then $\omega^{-1}(\psi(y))$ is, as one sees at once, the measure function of $\varphi(\omega(x))$, and
\[
 f\bigl(\omega[\omega^{-1}(\psi(y))]\bigr)=f(\psi(y)).
\]

This now permits us to extend the definition of $\int f(x)\,d\varphi(x)$ to the case in which $\varphi(x)$ is defined for every real number and the integration is correspondingly from $-\infty$ to $+\infty$ (we again suppose that $f(x)$ is bounded and $\Phi$-measurable, and that $\varphi(x)$ is bounded, monotonically nondecreasing, and right-continuous). Namely, let
\[
 \xi=\varrho(x),\qquad x=\varrho^{-1}(\xi),
\]
be a continuous, monotonically increasing, and hence one-to-one map of $0<x<\alpha$ onto $-\infty<\xi<\infty$. We put
\begin{equation}
 \int_{-\infty}^{\infty}f(\xi)\,d\varphi(\xi)
 =\int_0^\alpha f(\varrho(x))\,d\varphi(\varrho(x)).
 \label{eq:LS-definition-line}
\end{equation}
The right-hand side defining this expression is independent of the choice of $\varrho(x)$. To see this for two choices $\varrho_1(x)$ and $\varrho_2(x)$, it suffices to put $\omega(x)=\varrho_2^{-1}(\varrho_1(x))$ in \hyperref[eq:LS-change-finite]{the finite-interval change-of-variable formula}. Properties \hyperref[prop:I.4.a]{a)}--\hyperref[prop:I.4.g]{g)} also hold for this integral, as is immediate. (In \hyperref[prop:I.4.e]{e)}--\hyperref[prop:I.4.g]{g)}, the condition on $\mathfrak M$ is to be stated in the form given in \noteref{21}, since we have not defined the ``measure function'' for a general $\varphi(x)$.) Moreover, the changes of variable $\xi=\varrho(x)$, $x=\varrho^{-1}(\xi)$ leave this integral unchanged as well, directly by the definition, provided only that all $x$ are mapped onto all $\xi$.

With the aid of $\int$, more precisely of $\int_{-\infty}^{\infty}$, we can now also define every $\int_{\overline{\mathfrak M}}$\noteref{22}. Given $f(x)$, put $f_{\overline{\mathfrak M}}(x)=f(x)$ on $\overline{\mathfrak M}$ and $f_{\overline{\mathfrak M}}(x)=0$ on its complement, and define
\begin{equation}
 \int_{\overline{\mathfrak M}}f(x)\,d\varphi(x)
 =\int_{-\infty}^{\infty}f_{\overline{\mathfrak M}}(x)\,d\varphi(x).
 \label{eq:LS-restricted}
\end{equation}
The validity of \hyperref[prop:I.4.a]{a)}--\hyperref[prop:I.4.g]{g)} for $\int$ immediately implies their validity for $\int_{\overline{\mathfrak M}}$. The definition of the latter, together with \hyperref[prop:I.4.b]{b)} and \hyperref[prop:I.4.g]{g)}, also gives at once:

\begin{description}
 \item[\textup{h)}]\phantomsection\label{prop:I.4.h}
 Let $\overline{\mathfrak M}_1,\overline{\mathfrak M}_2,\ldots$ be pairwise disjoint $\Phi$-measurable sets (finitely or countably many), and let $\overline{\mathfrak M}=\overline{\mathfrak M}_1\cup\overline{\mathfrak M}_2\cup\cdots$. Then
\end{description}

% Source PDF page 10 / journal p. 200
\begin{equation}
 \int_{\overline{\mathfrak M}_1}f(x)\,d\varphi(x)
 +\int_{\overline{\mathfrak M}_2}f(x)\,d\varphi(x)+\cdots
 =\int_{\overline{\mathfrak M}}f(x)\,d\varphi(x).
 \label{eq:LS-countable-additivity}
\end{equation}

Finally, when $\overline{\mathfrak M}$ is the set $0\leq x\leq\alpha$, the integral $\int_{\overline{\mathfrak M}}$ is identical with the integral over the interval $0,\alpha$ defined at the beginning. To see this, define $f(x)=0$ outside $0\leq x\leq\alpha$, and define $\varphi(x)$ there as $\varphi(0)$ on the left and $\varphi(\alpha)$ on the right. Consider
\[
 \int_0^{\alpha+2\varepsilon}f(x-\varepsilon)\,d\varphi(x-\varepsilon),
 \qquad \varepsilon>0.
\]
If $\psi(y)$ is the measure function of $\varphi(x)$ on $0,\alpha$, then the measure function of $\varphi(x-\varepsilon)$ on $0,\alpha+2\varepsilon$ equals $\psi(y)+\varepsilon$ for $0\leq y<\varphi(\alpha)$, and equals $\alpha+2\varepsilon$ (not $\psi(y)+\varepsilon=\alpha+\varepsilon$) for $y=\varphi(\alpha)$; indeed, $0\leq y\leq\varphi(\alpha)$ is the domain. It follows that
\begin{equation}
 \int_0^{\alpha+2\varepsilon}f(x-\varepsilon)\,d\varphi(x-\varepsilon)
 =\int_0^\alpha f(x)\,d\varphi(x).
 \label{eq:LS-finite-line-agreement}
\end{equation}
The left-hand side is, however, equal to $\int_{-\infty}^{\infty}f(x)\,d\varphi(x)$. For this it suffices to choose the map $\varrho(x)$ from $0,\alpha+2\varepsilon$ onto $-\infty,+\infty$ so that $\varrho(x)=x-\varepsilon$ for $\varepsilon<x<\alpha+\varepsilon$, which is possible. Since $f(x)$ vanishes outside $0\leq x\leq\alpha$, this last integral is equal to $\int_{\overline{\mathfrak M}}f(x)\,d\varphi(x)$ when $\overline{\mathfrak M}$ is that set.

It is also worth recording that, by the definition, the change-of-variable formula takes the form
\begin{equation}
 \int_{\overline{\mathfrak M}}f(\xi)\,d\varphi(\xi)
 =\int_{\varrho^{-1}(\overline{\mathfrak M})}
 f(\varrho(x))\,d\varphi(\varrho(x)),
 \label{eq:LS-change-restricted}
\end{equation}
where $\varrho^{-1}(\overline{\mathfrak M})$ is the image of $\overline{\mathfrak M}$ under $x=\varrho^{-1}(\xi)$. Furthermore, in \hyperref[prop:I.4.e]{e)}--\hyperref[prop:I.4.g]{g)} (cf. \noteref{21}), the complement of $\mathfrak M$ may be replaced by the part of that complement lying in $\overline{\mathfrak M}$: outside $\overline{\mathfrak M}$ one has in any event
\[
 f_{\overline{\mathfrak M}}(x)
 =g_{\overline{\mathfrak M}}(x)
 =f_{n,\overline{\mathfrak M}}(x)=0,
\]
so all the requirements are fulfilled there.

\subsection*{5.}
\label{sec:I.5}
We now proceed to define $\int_{\overline{\mathfrak M}}f(x)\,d\varphi(x)$ for functions $\varphi(x)$ of bounded variation on the whole interval $-\infty<x<+\infty$. (As before, $f(x)$ is to be bounded and $\Phi$-measurable, and $\overline{\mathfrak M}$ is to be $\Phi$-measurable.) Such a function $\varphi(x)$ can, as is well known, always be put in the form
\begin{equation}
 \varphi(x)=\psi_1(x)-\psi_2(x)
 +i\bigl(\psi_3(x)-\psi_4(x)\bigr),
 \label{eq:BV-decomposition}
\end{equation}
where $\psi_1,\psi_2,\psi_3,\psi_4$ are $M$-functions (that is, real, never negative--although $\varphi(x)$ itself may be complex--monotonically nondecreasing, and right-continuous); this equality holds

% Source PDF page 11 / journal p. 201
except at the finitely or countably many points $x$ at which $\varphi(x)$ is not right-continuous\noteref{23}. We now define
\begin{align}
 \int_{\overline{\mathfrak M}}f(x)\,d\varphi(x)
 &={}
 \int_{\overline{\mathfrak M}}f(x)\,d\psi_1(x)
 -\int_{\overline{\mathfrak M}}f(x)\,d\psi_2(x)\notag\\
 &\quad
 +i\left(
 \int_{\overline{\mathfrak M}}f(x)\,d\psi_3(x)
 -\int_{\overline{\mathfrak M}}f(x)\,d\psi_4(x)
 \right).
 \label{eq:LS-BV-definition}
\end{align}
This is independent of the particular choice of $\psi_1,\psi_2,\psi_3,\psi_4$. Indeed,
\[
 \int_{\overline{\mathfrak M}}f(x)\,d\chi_1(x)
 -\int_{\overline{\mathfrak M}}f(x)\,d\chi_2(x),
\]
for $M$-functions $\chi_1,\chi_2$, depends only on $\chi_1(x)-\chi_2(x)$, as follows immediately from \hyperref[prop:I.4.d]{d)}. For an $M$-function $\varphi(x)$ the new definition agrees with the old one, since one may then take $\psi_1=\varphi$ and $\psi_2=\psi_3=\psi_4=0$.

Properties \hyperref[prop:I.4.a]{a)}--\hyperref[prop:I.4.h]{h)} carry over to this integral, as follows directly from the definition, and so does the change-of-variable formula at the end of \hyperref[sec:I.4]{subsection 4}. The following qualifications should nevertheless be added. For \hyperref[prop:I.4.c]{c)}, the assertion now holds for every complex $a$. It is enough to prove this for $M$-functions $\varphi(x)$; for them it already holds for positive $a$, and, by \hyperref[prop:I.4.d]{d)}, for $a=\pm1,\pm i$. These numbers are powers of $i$, so it suffices to take $a=i$, and this is immediate from the definition. For \hyperref[prop:I.4.e]{e)}--\hyperref[prop:I.4.g]{g)}, use the condition on $\mathfrak M$ formulated in \noteref{21}; it must be applied here to $\psi_1,\psi_2,\psi_3,\psi_4$, not to $\varphi$ itself. Note, however, that when the $\psi_1,\psi_2,\psi_3,\psi_4$ are chosen as in \noteref{23}, they are discontinuous only where $\varphi$ is discontinuous. For \hyperref[prop:I.4.e]{e)}, the assertions with $\geq0$ or $\leq0$ hold only for $M$-functions $\varphi(x)$, whereas the assertion with equality to zero plainly holds in general. If $f(x)$ and $\varphi(x)$ are real, the integral is plainly real as well.

Finally we prove the following important relation:
\begin{description}
 \item[\textup{i)}]\phantomsection\label{prop:I.5.i}
 Let $f(x)$ and $g(x)$ be bounded and $\Phi$-measurable, and let $\varphi(x)$ be of bounded variation. Then, as a function of $x$,
 \[
  \int^x g(y)\,d\varphi(y)
 \]
 is also of bounded variation (here $\int^x$ denotes integration over the set of $y\leq x$), and
 \begin{equation}
  \int f(x)\,d\left[\int^x g(y)\,d\varphi(y)\right]
  =\int f(x)g(x)\,d\varphi(x).
  \label{eq:LS-substitution-nested}
 \end{equation}
\end{description}

Since each of $f(x)$ and $g(x)$ can be expressed additively, with the factors $\pm1,\pm i$, in terms of nonnegative (and still bounded and $\Phi$-measurable) functions, and likewise every $\varphi(x)$ in terms of nonnegative, bounded, monotonically nondecreasing functions, by \hyperref[prop:I.4.a]{a)}--\hyperref[prop:I.4.d]{d)}, we may restrict ourselves in advance

% Source PDF page 12 / journal p. 202
to such functions. Then
\[
 \int^x g(y)\,d\varphi(y)
\]
is monotonically nondecreasing and bounded,
\[
 0\leq\int^x g(y)\,d\varphi(y)
 \leq\int g(y)\,d\varphi(y),
\]
by \hyperref[prop:I.4.e]{e)} and \hyperref[prop:I.4.h]{h)}, and is therefore of bounded variation. This proves the first assertion. It remains to prove the integral formula. By \hyperref[prop:I.4.f]{f)} and \hyperref[thm:3]{Theorem 3}, we need consider only $f(x)$ and $g(x)$ belonging to the third Baire class. By the ``continuity'' in $f(x)$ of both sides expressed in \hyperref[prop:I.4.g]{g)}, we may even restrict $f(x)$ to the zeroth Baire class, that is, assume it continuous. For continuous $f(x)$, however, the left-hand side
\[
 \int f(x)\,d\eta(x),
 \qquad
 \eta(x)=\int^x g(y)\,d\varphi(y),
\]
is an ordinary Riemann--Stieltjes integral and is therefore ``continuous'' in $\eta(x)$. By \hyperref[prop:I.4.g]{g)}, both $\eta(x)$ and the right-hand side $\int f(x)g(x)\,d\varphi(x)$ are ``continuous'' in $g(x)$, so we may also regard $g(x)$ as belonging to the zeroth Baire class, that is, as continuous. Both sides then involve only Riemann--Stieltjes integrals and are therefore also ``continuous'' in $\varphi(x)$. Since there certainly exists a sequence of everywhere continuous and differentiable (nonnegative and monotonically nondecreasing) functions converging everywhere to $\varphi(x)$, we may suppose that $\varphi(x)$ itself has these properties. We may then write the Riemann--Stieltjes integrals as Riemann integrals:
\begin{align}
 \int f(x)\,d\left[\int^x g(y)\,d\varphi(y)\right]
 &=\int f(x)\,
 \frac{d}{dx}\left[\int^x g(y)\frac{d\varphi(y)}{dy}\,dy\right]dx\notag\\
 &=\int f(x)g(x)\frac{d\varphi(x)}{dx}\,dx,
 \label{eq:nested-RS-left}
\\
 \int f(x)g(x)\,d\varphi(x)
 &=\int f(x)g(x)\frac{d\varphi(x)}{dx}\,dx.
 \label{eq:nested-RS-right}
\end{align}
This makes the asserted relation evident.

\section*{II. General Functions of Operators}
\label{sec:II}

\subsection*{1.}
\label{sec:II.1}
Let $A$ be a bounded Hermitian operator, hence an element of $\mathbf B$. We write $A$ in Hilbert's spectral form\noteref{24}:
\begin{equation}
 (Af,g)=\int \lambda\,d(E(\lambda)f,g),
 \label{eq:Hilbert-spectral-form}
\end{equation}
where $\lambda$ is the variable of integration and $E(\lambda)$ is the resolution of the identity belonging to $A$. (The notation is that explained in E and A at the places cited in \noteref{24}.) By $f,g,\ldots$ and $\varphi,\psi,\ldots$ we shall henceforth understand

% Source PDF page 13 / journal p. 203
vectors of the Hilbert space (general functions [defined on the real line] will be denoted by
\(F,G,\ldots\)).  The integral is a Riemann--Stieltjes integral\noteref{25}, and
it is certainly meaningful (despite the unbounded \(\lambda\)!), since
\(E(\lambda)\), and hence also \((E(\lambda)f,g)\), is constant outside a finite
interval\noteref{26}; it therefore suffices to integrate over that interval.

Consider \(E(\lambda)\) as a function of \(\lambda\).  With respect to the
relation \(\leq\) defined for projection operators, it is monotone
nondecreasing and right-continuous\noteref{27}; below a suitable bound it
equals \(0\), and above a suitable bound it equals \(1\).  Let \(\mathfrak T\)
be the union of its intervals of constancy (regarded as open), and let
\(\mathfrak S\) be the complement of \(\mathfrak T\).  Thus \(\mathfrak T\) is
open, \(\mathfrak S\) is closed, and, by what was just said, \(\mathfrak S\) is
bounded.  For every \(f,g\), the function \((E(\lambda)f,g)\) is constant on
the intervals of constancy of \(E(\lambda)\); that is, \(\mathfrak T\) is
covered by intervals of constancy of this function.  As is well known,
\(\mathfrak S\) is the ``spectrum'' of \(A\).  Further, let \(\mathfrak P\) be
the set of all discontinuity points of \(E(\lambda)\).  Only at these points
can a function \((E(\lambda)f,g)\) be discontinuous (cf.\ \noteref{27}).
Naturally, \(\mathfrak P\subseteq\mathfrak S\), and \(\mathfrak P\) is known
to be finite or countable; it is the ``point spectrum'' of \(A\).

Now let \(F(x)\) be a bounded \(\Phi\)-measurable function.  We consider the
expressions
\[
  \int F(\lambda)\,d(E(\lambda)f,g).
\]
Since \(\mathfrak T\) consists entirely of intervals of constancy of the
function following the differential sign, all these expressions are
independent of the behavior of \(F(x)\) on \(\mathfrak T\)
(\hyperref[sec:I.4]{I.4 and I.5, property (f)}, together with
\noteref{21}).  We may therefore, for example, put \(F(x)=0\) on
\(\mathfrak T\) (since \(\mathfrak T\) is open, \(F(x)\) remains
\(\Phi\)-measurable), or, what is the same thing, replace the integral by an
integral over \(\mathfrak S\).  Consequently \(F(x)\) need not be bounded
everywhere; it is enough that it be bounded on \(\mathfrak S\)\noteref{28}.

In general, the following estimate holds:
\[
 \left|\int F(\lambda)\,d(E(\lambda)f,g)\right|
 \leq
 \sqrt{
   \left(\int |F(\lambda)|\,d|E(\lambda)f|^2\right)
   \left(\int |F(\lambda)|\,d|E(\lambda)g|^2\right)}.
\]
It suffices first to prove this for functions \(F(\lambda)\) of the third
Baire class (\hyperref[thm:3]{Theorem~3} and
\hyperref[sec:I.4]{I.4(f)}), and consequently even for those of the zeroth
class (\hyperref[sec:I.4]{I.4(g)}), that is, for continuous functions.  The
integrals are then Riemann--Stieltjes integrals and hence, respectively,
limits of
\begin{align*}
 &\sum_{n=1}^{N} F(\lambda_n)
 \bigl\{(E(\lambda_n)f,g)-(E(\lambda_{n-1})f,g)\bigr\} \\
 &\qquad={}
 \sum_{n=1}^{N}F(\lambda_n)
 \bigl(\{E(\lambda_n)-E(\lambda_{n-1})\}f,g\bigr),
\end{align*}

% Source PDF page 14 / journal p. 204
\begin{align*}
 &\sum_{n=1}^{N}|F(\lambda_n)|
 \bigl\{|E(\lambda_n)f|^2-|E(\lambda_{n-1})f|^2\bigr\} \\
 &\qquad={}
 \sum_{n=1}^{N}|F(\lambda_n)|
 \bigl|\{E(\lambda_n)-E(\lambda_{n-1})\}f\bigr|^2,\noteref{29}\\
 &\sum_{n=1}^{N}|F(\lambda_n)|
 \bigl\{|E(\lambda_n)g|^2-|E(\lambda_{n-1})g|^2\bigr\} \\
 &\qquad={}
 \sum_{n=1}^{N}|F(\lambda_n)|
 \bigl|\{E(\lambda_n)-E(\lambda_{n-1})\}g\bigr|^2,\noteref{29}
\end{align*}
where
\begin{align*}
 \lambda_0&<\lambda_1<\cdots<\lambda_N,\\
 \lambda_0&\longrightarrow-\infty,
 &\lambda_N&\longrightarrow+\infty,\\
 \max_{n=1,\ldots,N}(\lambda_n-\lambda_{n-1})&\longrightarrow0.
\end{align*}
Because
\[
 \left|\bigl(\{E(\lambda_n)-E(\lambda_{n-1})\}f,g\bigr)\right|
 \leq
 \bigl|\{E(\lambda_n)-E(\lambda_{n-1})\}f\bigr|
 \bigl|\{E(\lambda_n)-E(\lambda_{n-1})\}g\bigr|,\noteref{30}
\]
and because of Schwarz's inequality
\[
 \left|\sum_{n=1}^{N}a_nb_n\right|
 \leq
 \sqrt{\left(\sum_{n=1}^{N}|a_n|^2\right)
       \left(\sum_{n=1}^{N}|b_n|^2\right)},
\]
the absolute value of the first expression is at most the square root of
the product of the latter two.  This proves the assertion.

Suppose now that, in general, \(|F(x)|\leq C\) (we know, incidentally, that
it is enough for this to hold on \(\mathfrak S\)).  Then
\begin{align*}
 \int |F(x)|\,d|E(\lambda)f|^2
 &\leq \int C\,d|E(\lambda)f|^2
  =C\bigl(|1f|^2-|0f|^2\bigr)=C|f|^2,\\
 \int |F(x)|\,d|E(\lambda)g|^2
 &\leq \int C\,d|E(\lambda)g|^2
  =C\bigl(|1g|^2-|0g|^2\bigr)=C|g|^2,
\end{align*}
and hence
\[
 \left|\int F(x)\,d(E(\lambda)f,g)\right|
 \leq C|f|\,|g|.
\]
Keeping \(f\) fixed and viewing
\[
 \int F(x)\,d(E(\lambda)f,g)=L(g)
\]
as a function of \(g\), we therefore have
\[
 L(a_1g_1+\cdots+a_ng_n)
 =\overline{a_1}L(g_1)+\cdots+\overline{a_n}L(g_n),
 \qquad |L(g)|\leq D|g|,
 \qquad D=C|f|.
\]
By a theorem of F. Riesz\noteref{31}, there is therefore one and only one
\(f^*\) such that \((f^*,g)=L(g)\) always holds.  We call this vector
\(A'f\) (it does, after all, depend on \(f\)); thus, by definition,
\[
 (A'f,g)=\int F(\lambda)\,d(E(\lambda)f,g).
\]

% Source PDF page 15 / journal p. 205
As one sees at once, \(A'\) is a linear operator, and
\(|(A'f,g)|\leq C|f|\,|g|\) implies \(|A'f|\leq C|f|\) (put
\(g=A'f\)); hence it is bounded as well, that is, \(A'\in \mathbf B\).
We shall now denote this operator \(A'\) by \(F(A)\).

\subsection*{2.}\phantomsection\label{sec:II.2}

We list the simplest properties of \(F(A)\):
\begin{description}
 \item[\(\mathrm{a})\)] \phantomsection\label{prop:II.2.a}For \(F(x)\equiv0\), \(F(x)\equiv1\), and
 \(F(x)\equiv x\), respectively, one has \(F(A)=0\), \(F(A)=1\), and
 \(F(A)=A\), respectively.
 \item[\(\mathrm{b})\)] \phantomsection\label{prop:II.2.b}If \(F(x)\equiv\overline{G(x)}\) (complex
 conjugate), then \(F(A)=G(A)^*\).\noteref{32}
 \item[\(\mathrm{c})\)] \phantomsection\label{prop:II.2.c}If \(F(x)\equiv aG(x)\), then \(F(A)=aG(A)\).
 \item[\(\mathrm{d})\)] \phantomsection\label{prop:II.2.d}If \(F(x)\equiv G(x)+H(x)\), then
 \(F(A)=G(A)+H(A)\).
 \item[\(\mathrm{e})\)] \phantomsection\label{prop:II.2.e}If \(F(x)\equiv G(x)H(x)\), then
 \(F(A)=G(A)H(A)\).
 \item[\(\mathrm{f})\)] \phantomsection\label{prop:II.2.f}If \(F(x)\equiv G(H(x))\), then
 \(F(A)=G(H(A))\).\noteref{33}
 \item[\(\mathrm{g})\)] \phantomsection\label{prop:II.2.g}The equality \(F(x)=G(x)\) on a set
 \(\mathfrak M\) implies \(F(A)=G(A)\), provided that
 \(\mathfrak M\) has the following property.  Let
 \(\xi=\rho(x)\), \(x=\rho^{-1}(\xi)\), be a monotone increasing,
 continuous, and therefore one-to-one, mapping of an interval
 \(0<x<\alpha\) onto \(-\infty<\xi<\infty\).  The particular choice of
 \(\rho(x)\) is immaterial, for, as is easily seen, the expression
 \(\rho(\Psi_f(y))\) to be defined is only apparently dependent on it.
 Indeed, \(\rho(\Psi_f(y))\) is to be understood as the inverse, in the
 sense sketched at the end of \hyperref[sec:I.1]{I.1}, of
 \(|E(\xi)f|^2\).  Form
 \[
   \Phi_f(x)=|E(\rho(x))f|^2
 \]
 and its measure function \(\Psi_f(y)\), and then form
 \(\rho(\Psi_f(y))\).  For every \(f\), the set of all \(y\) for which
 \(\rho(\Psi_f(y))\notin\mathfrak M\) must be a Lebesgue null
 set.\noteref{34}
 \item[\(\mathrm{h})\)] \phantomsection\label{prop:II.2.h}Suppose \(F_n(x)\to G(x)\) as \(n\to\infty\)
 on a set \(\mathfrak M\) as in \(g\), and suppose further that the
 functions \(F_n(x)\) are uniformly bounded (at least on a set
 \(\mathfrak M\) of the preceding kind).  Then \(F_n(A)\to G(A)\) in
the sense of ``strong double convergence'' in \(\mathbf B\)\noteref{35};
 that is,
 \[
   F_n(A)f\longrightarrow G(A)f,
   \qquad F_n(A)^*f\longrightarrow G(A)^*f
 \]
 strongly in the Hilbert space.
\end{description}

Assertions \hyperref[prop:II.2.a]{(a)}--\hyperref[prop:II.2.h]{(h)}
say that a number of properties of the real numbers
\(x\) (on which alone our functions are directly defined) carry over to
all bounded Hermitian operators.  We shall now prove them.

Assertions \hyperref[prop:II.2.a]{(a)}--\hyperref[prop:II.2.d]{(d)} and
\hyperref[prop:II.2.g]{(g)} follow immediately from the definition of
\(F(A)\), and so forth (that is, from the definitions of the scalar
products \((F(A)f,g)\), and so forth).  In
\hyperref[prop:II.2.g]{(g)}, one must observe that the
condition imposed there on \(\mathfrak M\) has precisely the right form:
the functions of bounded variation that follow the differential sign in
the relevant definitions, namely \((E(\lambda)f,g)\), are to be regarded
as linear combinations of four functions \(|E(\lambda)h|^2\), with
\[
 h=\frac{f\pm g}{2},\qquad h=\frac{f\pm ig}{2},
\]
and with the corresponding coefficients \(\pm1,\pm i\) (cf.\
\noteref{25}).  Moreover, because the integrals extend over the interval
\(-\infty<\lambda<\infty\), they must then be reduced to an interval
\(0<x<\alpha\) by \(\lambda=\rho(x)\), \(x=\rho^{-1}(\lambda)\)
(cf.\ \hyperref[sec:I.4]{I.4}); this leads exactly to our condition on
\(\mathfrak M\).

% Source PDF page 16 / journal p. 206
One further point is worth mentioning.  Since both sides of the equality
to be proved,
\[
 (F(A)f,g)=(G(A)f,g),
\]
are continuous in \(f\), it is enough to prove it for \(f,g\) in an
everywhere dense set; because the Hilbert space is separable\noteref{36},
this set may be chosen countable.  The values of \(h\) that occur above
(called \(f\) in \hyperref[prop:II.2.g]{(g)}), namely
\((f\pm g)/2\) and \((f\pm ig)/2\), then
also form a countable set.  Thus it suffices to impose the requirements on
\(\mathfrak M\) in \hyperref[prop:II.2.g]{(g)} for a suitable countable
set of vectors \(f\).

Likewise, \hyperref[prop:II.2.h]{(h)} follows immediately if
\(F_n(A)\to F(A)\) there is asserted
only in the sense of ``weak'' convergence in \(\mathbf B\)
(cf.\ \noteref{35}); that is, if one need prove only
\[
 (F_n(A)f,g)\longrightarrow(F(A)f,g)
 \qquad(n\to\infty).
\]
To carry over also the preceding supplementary observation concerning
\hyperref[prop:II.2.g]{(g)}
(namely, the sufficiency of considering a suitable countable set of
vectors \(f\) in the condition on \(\mathfrak M\)), it is enough to recall
that the \(F_n(x)\), \(n=1,2,\ldots\), are uniformly bounded functions of
\(x\).  Hence the \(F_n(A)\) are uniformly bounded operators
(cf.\ \hyperref[sec:II.1]{II.1}), and the \((F_n(A)f,g)\) are therefore
equicontinuous functions of \(f,g\).

For \hyperref[prop:II.2.e]{(e)}, one has to prove, for all \(f,g\), the
equality of
\[
 (F(A)f,g)
 =\int F(\lambda)\,d(E(\lambda)f,g)
 =\int G(\lambda)H(\lambda)\,d(E(\lambda)f,g)
\]
and
\begin{align*}
 (G(A)H(A)f,g)
 &=\bigl(H(A)f,G(A)^*g\bigr)\\
 &=\int H(\lambda)\,d\bigl(E(\lambda)f,G(A)^*g\bigr)\\
 &=\int H(\lambda)\,d\bigl(G(A)E(\lambda)f,g\bigr).
\end{align*}
Now
\begin{align*}
 \bigl(G(A)E(\lambda)f,g\bigr)
 &=\int G(\lambda')\,d\bigl(E(\lambda')E(\lambda)f,g\bigr)\\
 &=\int G(\lambda')\,d\bigl(E(\min(\lambda',\lambda))f,g\bigr).
\end{align*}
The expression following the differential sign is constant for
\(\lambda'\geq\lambda\) (for then
\(\min(\lambda',\lambda)=\lambda\); \(\lambda\) is fixed and
\(\lambda'\) is the variable).  We may therefore replace the integral by
\(\int_{-\infty}^{\lambda}\) (cf.\
\hyperref[sec:I.4]{I.4(e)} and \noteref{21}); throughout this range of
integration \(\lambda'\leq\lambda\), so
\(\min(\lambda',\lambda)=\lambda'\).  Hence
\[
 \bigl(G(A)E(\lambda)f,g\bigr)
 =\int_{-\infty}^{\lambda}G(\lambda')\,
   d(E(\lambda')f,g).
\]
It remains to prove
\[
 \int H(\lambda)G(\lambda)\,d(E(\lambda)f,g)
 =\int H(\lambda)\,d\left[
    \int_{-\infty}^{\lambda}G(\lambda')\,d(E(\lambda')f,g)
  \right],
\]
and this is precisely the assertion of \hyperref[sec:I.4]{I.4(i)}.

% Source PDF page 17 / journal p. 207
We can now complete the proof of \hyperref[prop:II.2.h]{(h)}.  Alongside
the sequence \(F_n(x)\)
and the function \(F(x)\), consider the sequence
\(\overline{F_n(x)}F_n(x)\) and the function
\(\overline{F(x)}F(x)\).  To both,
\hyperref[prop:II.2.h]{(h)} applies to the extent already
proved.  Thus \(F_n(A)\) converges weakly in \(\mathbf B\) to \(F(A)\),
and \(F_n(A)^*F_n(A)\) likewise converges weakly to \(F(A)^*F(A)\)
(here one uses \hyperref[prop:II.2.b]{(b)} and
\hyperref[prop:II.2.e]{(e)}, just proved).  In other words, for all \(f,g\),
\begin{align*}
 (F_n(A)f,g)&\longrightarrow(F(A)f,g),\\
 (F_n(A)^*F_n(A)f,g)&\longrightarrow(F(A)^*F(A)f,g),
\end{align*}
and hence
\[
 (F_n(A)f,F_n(A)g)\longrightarrow(F(A)f,F(A)g).
\]
Taking \(f=g\), the last assertion says
\(|F_n(A)f|\to|F(A)f|\).  As is well known, it now follows that
\(F_n(A)f\to F(A)f\) in the strong topology of the Hilbert
space\noteref{37}.  Finally, the ``double convergence'' follows on applying
the same argument to \(\overline{F_n(x)}\) and \(\overline{F(x)}\).

It remains to prove \hyperref[prop:II.2.f]{(f)}.  By
\hyperref[prop:II.2.g]{(g)}, we may alter the functions \(G(x)\) and
\(H(x)\) occurring there (and with them \(F(x)=G(H(x))\)) on certain sets
without affecting the result, provided these sets have the following
properties.  Leave \(H(x)\) unchanged, and alter \(G(x)\) on a set
\(\mathfrak A\).  In order that \(G(H(A))\) remain invariant, for each
\(f\) the set of those \(y\) for which \(\rho(\Psi_f(y))\), formed as in
\hyperref[prop:II.2.g]{(g)} for \(H(A)\), belongs to \(\mathfrak A\),
must be a Lebesgue null set.
In order that \(F(A)\) remain invariant, for each \(f\) the set of those
\(y\) for which \(\rho(\Psi_f(y))\), now formed as in
\hyperref[prop:II.2.g]{(g)} for \(A\),
belongs to the set on which \(F(x)\) is altered--that is, for which
\(H(\rho(\Psi_f(y)))\in\mathfrak A\)--must be a Lebesgue null set.  By
the observation following the proof of \hyperref[prop:II.2.g]{(g)}, it is
enough in each case to
consider a suitable countable set of vectors \(f\).  Consequently
\hyperref[thm:3]{Theorem~3} applies: we may suppose that \(G(x)\) belongs
to the third Baire class, and, by \hyperref[prop:II.2.h]{(h)}, even to the
zeroth class, which in
turn may be taken to consist of the polynomials.  Thus let \(G(x)\) be a
polynomial.

If
\[
 F(x)=a_0(H(x))^n+a_1(H(x))^{n-1}+\cdots+a_n,
\]
then, by \hyperref[prop:II.2.a]{(a)}--\hyperref[prop:II.2.d]{(d)},
\[
 F(A)=a_0(H(A))^n+a_1(H(A))^{n-1}+\cdots+a_n\,1.
\]
On the other hand, putting \(H(x)=x\), for
\[
 G(x)=a_0x^n+a_1x^{n-1}+\cdots+a_n
\]
we obtain
\[
 G(B)=a_0B^n+a_1B^{n-1}+\cdots+a_n\,1
 \qquad\text{(cf. (a)).}
\]
Thus \hyperref[prop:II.2.f]{(f)} holds for every polynomial \(G(x)\).

\subsection*{3.}\phantomsection\label{sec:II.3}

We are now in a position to prove the following theorem on the convergence
of sequences of operator functions.

\medskip
\phantomsection\label{thm:4}\noindent\textbf{Theorem 4.}
\emph{Let \(F_1(x),F_2(x),\ldots\) be a sequence of
\(\Phi\)-measurable functions, and let \(A\) be a Hermitian operator.
Then the following assertions are equivalent:}
\begin{description}
 \item[\(\alpha)\)] \phantomsection\label{thm:4.alpha}\emph{For every \(f\in\mathfrak H\), the vectors
 \(F_1(A)f,F_2(A)f,\ldots\) converge in \(\mathfrak H\) to a limit, in
 the sense of the strong topology of \(\mathfrak H\)
 (cf.\ \noteref{35}).}

% Source PDF page 18 / journal p. 208
 \item[\(\beta)\)] \phantomsection\label{thm:4.beta}\emph{Assertion
 \hyperref[thm:4.alpha]{\(\alpha\)} holds both for
 \(F_1(A)f,F_2(A)f,\ldots\) and for
 \(F_1(A)^*f,F_2(A)^*f,\ldots\).}
 \item[\(\gamma)\)] \phantomsection\label{thm:4.gamma}\emph{All the functions
 \(F_1(x),F_2(x),\ldots\) are uniformly bounded on a set
 \(\mathfrak M\) having the property specified in
 \hyperref[prop:II.2.g]{II.2(g)}.  If \(m_1,m_2,\ldots\) and
 \(n_1,n_2,\ldots\) are two sequences of integers tending to infinity,
 and \(f\) is an arbitrary element of \(\mathfrak H\), then the sequence
 \[
  F_{m_1}(x)-F_{n_1}(x),\quad
  F_{m_2}(x)-F_{n_2}(x),\quad\ldots
 \]
 has a subsequence
 \[
  F_{m_{p_1}}(x)-F_{n_{p_1}}(x),\quad
  F_{m_{p_2}}(x)-F_{n_{p_2}}(x),\quad\ldots,
  \qquad p_1<p_2<\cdots,
 \]
 such that, for this \(f\), the set \(\mathfrak M\) of all \(x\) on
 which it converges to \(0\) has the property specified in
 \hyperref[prop:II.2.g]{II.2(g)}.}
\end{description}

\noindent\textit{Proof.}
Assertion \hyperref[thm:4.beta]{\(\beta\)} is obtained from
\hyperref[thm:4.alpha]{\(\alpha\)} by adjoining the same
requirement for \(F_1(A)^*,F_2(A)^*,\ldots\), that is, for
\(\overline{F_1(x)},\overline{F_2(x)},\ldots\).  Since
\hyperref[thm:4.gamma]{\(\gamma\)} has
exactly the same form for these functions, it is enough to compare
\hyperref[thm:4.alpha]{\(\alpha\)} and
\hyperref[thm:4.gamma]{\(\gamma\)}.

First consider the case in which \hyperref[thm:4.gamma]{\(\gamma\)}
holds.  If \hyperref[thm:4.alpha]{\(\alpha\)} were
false, there would be an \(f\) for which the vectors \(F_n(A)f\) have no
strong limit in \(\mathfrak H\); equivalently, by the topological
completeness of \(\mathfrak H\), Cauchy's convergence criterion would
fail.  Thus there is an \(\epsilon>0\) such that, for arbitrarily large
\(m\), a suitable \(n\geq m\) exists for which
\[
 |(F_m(A)-F_n(A))f|
 =|F_m(A)f-F_n(A)f|\geq\epsilon.
\]
We can therefore find two sequences \(m_1,m_2,\ldots\) and
\(n_1,n_2,\ldots\), both tending to infinity, such that
\[
 |(F_{m_v}(A)-F_{n_v}(A))f|\geq\epsilon
 \qquad(v=1,2,\ldots).
\]
In particular, for the \(p_1,p_2,\ldots\) supplied by
\hyperref[thm:4.gamma]{\(\gamma\)}, it
cannot be the case that
\((F_{m_{p_v}}(A)-F_{n_{p_v}}(A))f\to0\), although this should follow
from \hyperref[prop:II.2.h]{II.2(h)}.\noteref{38}

Now consider the second case, in which
\hyperref[thm:4.alpha]{\(\alpha\)} holds.  We prove the first half of
\hyperref[thm:4.gamma]{\(\gamma\)} as follows.  Since \(F_n(A)f\)
converges for every
\(f\), the operators \(F_n(A)\) are uniformly bounded (cf.\
\hyperref[ref:A]{A}, p.~382; the Banach--Saks proof sketched there for
strongly convergent \(F_n(A)\) applies here without change).  Hence there
is a number \(C\) such that
\[
 |F_n(A)f|\leq C|f|
\]
for every \(f\) and \(n\).  We shall therefore have reached our goal if we
can prove the following general assertion: if
\(|G(A)f|\leq C|f|\) for every \(f\), then \(|G(x)|\leq C\) on a set
\(\mathfrak M\) having the property specified in
\hyperref[prop:II.2.g]{II.2(g)}.  Indeed, the intersection of the sets
\(\mathfrak M\) so obtained for \(F_1(x),F_2(x),\ldots\) clearly has that
property again.

For real \(G(x)\), and hence Hermitian \(G(A)\), known theorems of Hilbert
imply from \(|G(A)f|\leq C|f|\) for all \(f\) that the spectrum of \(G(A)\)
lies entirely in the interval \(-C\leq x\leq C\).  If a function \(H(x)\)
vanishes everywhere on this interval, then
\(H(G(A))=0\) by \hyperref[prop:II.2.g]{II.2(g)} and \noteref{34}.
Applying the concluding observation of \hyperref[sec:II]{II} to
\(H(G(x))\), it follows that \(H(G(x))=0\) on a set \(\mathfrak M\) of
the stated kind.  Put
\[
 H(x)=
 \begin{cases}
  0,&-C\leq x\leq C,\\
  1,&\text{otherwise}.
 \end{cases}
\]

% Source PDF page 19 / journal p. 209
This proves our assertion for real \(G(x)\).  For complex \(G(x)\), we
have \(|G(A)f|\leq C|f|\).  Put \(H(x)=\overline{G(x)}\); then
\[
 |H(A)f|=|G(A)^*f|\leq C|f|
\]
(cf.\ \hyperref[ref:E]{E}, Appendix II, p.~112).  Hence, for
\[
 K(x)=\Re G(x)=\frac12\bigl(G(x)+H(x)\bigr),
\]
we have \(|K(A)f|\leq C|f|\).  It follows that
\(|\Re G(x)|\leq C\) on a set \(\mathfrak M\).  Replacing \(G(x)\) by
\(\theta G(x)\), where \(\theta\) is constant and \(|\theta|=1\), gives
in exactly the same way
\[
 |\Re(\theta G(x))|\leq C
\]
on a set \(\mathfrak M\).  Choose a sequence of values of \(\theta\)
dense on \(|\theta|=1\).  The preceding inequality then holds
simultaneously for all \(\theta\) in this sequence on the intersection of
a sequence of such sets, and hence again on a set \(\mathfrak M\).
Wherever this is so, it must hold for every \(\theta\) with
\(|\theta|=1\), and then \(|G(x)|\leq C\).

It remains to prove the second half of
\hyperref[thm:4.gamma]{\(\gamma\)}.  If it were false,
there would be two sequences \(m_1,m_2,\ldots\) and \(n_1,n_2,\ldots\),
and an \(f\), such that no subsequence of
\[
 F_{m_1}(x)-F_{n_1}(x),\quad
 F_{m_2}(x)-F_{n_2}(x),\quad\ldots
\]
converges to \(0\) on a set containing all the
\(\rho(\Psi_f(y))\), apart from a Lebesgue null set of \(y\)'s.

Suppose that for every \(\epsilon>0\) and \(\eta>0\) there is a
\(v_0=v_0(\epsilon,\eta)\) such that, for every fixed \(v\geq v_0\), the
set\tnref{C}
\[
 \mathfrak M_v(\epsilon)
 =\bigl\{x:|F_{m_v}(x)-F_{n_v}(x)|<\epsilon\bigr\}
\]
contains all \(\rho(\Psi_f(y))\), with the exception of a set of \(y\)'s
of Lebesgue measure at most \(\eta\).  Choose
\begin{align*}
 p_1&\geq v_0\left(\frac12,\frac12\right),\\
 p_2&>p_1, &p_2&\geq v_0\left(\frac14,\frac14\right),\\
 p_3&>p_2, &p_3&\geq v_0\left(\frac18,\frac18\right),\quad\ldots.
\end{align*}
For every \(x\) in the intersection
\[
 \mathfrak M_{p_\mu}\left(\frac1{2^\mu}\right)
 \cap\mathfrak M_{p_{\mu+1}}\left(\frac1{2^{\mu+1}}\right)
 \cap\cdots,
\]
we then have
\[
 |F_{m_{p_v}}(x)-F_{n_{p_v}}(x)|<\frac1{2^v}
 \qquad\text{for every }v\geq\mu.
\]
This intersection contains all \(\rho(\Psi_f(y))\), apart from a set of
\(y\)'s of Lebesgue measure at most
\[
 \frac1{2^\mu}+\frac1{2^{\mu+1}}+\cdots=\frac1{2^{\mu-1}}.
\]
For these \(x\), the sequence
\(F_{m_{p_v}}(x)-F_{n_{p_v}}(x)\) plainly converges to \(0\).  If
\(\mathfrak M\) denotes the set of \(x\)'s on which it converges, all
\(\rho(\Psi_f(y))\) belong to \(\mathfrak M\), apart from a set of
\(y\)'s of Lebesgue measure at most \(2^{1-\mu}\).  This holds for every
\(\mu\), so that this exceptional set has Lebesgue measure \(0\), contrary
to what was established above.  There must therefore be
\(\epsilon>0\) and \(\eta>0\) such that, for infinitely many \(v\), the
set \(\mathfrak M_v(\epsilon)\) fails to contain the
\(\rho(\Psi_f(y))\) on a set of \(y\)'s of Lebesgue measure at least
\(\eta\).

For these \(v\), however,
\begin{align*}
 &|(F_{m_v}(A)-F_{n_v}(A))f|^2\\
 &\quad=
 \bigl((F_{m_v}(A)-F_{n_v}(A))^*
       (F_{m_v}(A)-F_{n_v}(A))f,f\bigr)\\
 &\quad=\int_{-\infty}^{\infty}
   |F_{m_v}(\xi)-F_{n_v}(\xi)|^2\,d|E(\xi)f|^2.
\end{align*}

% Source PDF page 20 / journal p. 210
By the changes of variable used above, this is
\begin{align*}
 &=\int_0^\alpha
   |F_{m_v}(\rho(x))-F_{n_v}(\rho(x))|^2\,
   d|E(\rho(x))f|^2\\
 &=\int_0^{|f|^2}
   \left|F_{m_v}(\rho(\Psi_f(y)))
        -F_{n_v}(\rho(\Psi_f(y)))\right|^2\,dy\\
 &\geq
 \int_{\substack{0\leq y\leq|f|^2\\
       \rho(\Psi_f(y))\notin\mathfrak M_v(\epsilon)}}
   \left|F_{m_v}(\rho(\Psi_f(y)))
        -F_{n_v}(\rho(\Psi_f(y)))\right|^2\,dy\\
 &\geq
 \int_{\substack{0\leq y\leq|f|^2\\
       \rho(\Psi_f(y))\notin\mathfrak M_v(\epsilon)}}
   \epsilon^2\,dy
 \geq\epsilon^2\eta.
\end{align*}
Consequently
\[
 \limsup_{v\to\infty}|F_{m_v}(A)f-F_{n_v}(A)f|
 \geq\sqrt{\epsilon^2\eta}>0,
\]
which is incompatible with the convergence of the \(F_n(A)f\), that is,
with \hyperref[thm:4.alpha]{\(\alpha\)}.  This proves the theorem.

We add the following observation concerning this theorem.  Assertions
\hyperref[thm:4.alpha]{\(\alpha\)} and
\hyperref[thm:4.beta]{\(\beta\)}, respectively, say that for every \(f\),
\(F_n(A)f\) has a strong limit in \(\mathfrak H\), and that both
\(F_n(A)f\) and \(F_n(A)^*f\) have such a limit.  The same \(f\) also
enters into assertion \hyperref[thm:4.gamma]{\(\gamma\)}.
\hyperref[thm:4]{Theorem~4} says that
the validity of \hyperref[thm:4.alpha]{\(\alpha\)} for all \(f\), that of
\hyperref[thm:4.beta]{\(\beta\)} for all \(f\), and that of
\hyperref[thm:4.gamma]{\(\gamma\)} for all \(f\), are equivalent.  The
proof, however, shows that \hyperref[thm:4.alpha]{\(\alpha\)} and
\hyperref[thm:4.beta]{\(\beta\)} follow from
\hyperref[thm:4.gamma]{\(\gamma\)} even
when all three are considered only for one fixed \(f\).  (Here and below,
the uniform boundedness of the \(F_n(x)\) on a set \(\mathfrak M\) having
the property of \hyperref[prop:II.2.g]{II.2(g)} for all \(f\) is assumed.)
Since the \(F_n(A)\) and \(F_n(A)^*\) are uniformly bounded (because the
\(F_n(x)\) are essentially uniformly bounded), the set of \(f\)'s in
\(\mathfrak H\) for which \hyperref[thm:4.alpha]{\(\alpha\)},
respectively \hyperref[thm:4.beta]{\(\beta\)}, holds is
a closed linear manifold.  Thus these assertions hold everywhere if, for
example, they hold for a complete normalized orthogonal system.  By the
equivalence just proved, the same is true of
\hyperref[thm:4.gamma]{\(\gamma\)}: it is not weakened
if its second half is required, rather than for every \(f\), only for the
(countably many) elements of any complete orthogonal system.

We next prove the following.

\medskip
\phantomsection\label{thm:5}\noindent\textbf{Theorem 5.}
\begin{itshape}
Let the \(F_n(x)\) and \(A\) be as in
\hyperref[thm:4]{Theorem~4}.  If the convergence described there takes
place, then there is an operator \(B\in \mathbf B\) such that
\[
 F_n(A)f\longrightarrow Bf,
 \qquad F_n(A)^*f\longrightarrow B^*f;
\]
that is, \(B\) is the strong double limit of the \(F_n(A)\)
(cf.\ \noteref{35}).  This \(B\) is again of the form \(F(A)\), where
\(F(x)\) may be chosen bounded and \(\Phi\)-measurable.

Moreover, the following condition is necessary and sufficient for \(F(A)\)
to be the strong (respectively, strong double) limit of the \(F_n(A)\): all
the \(F_n(x)\) are uniformly bounded on a set \(\mathfrak M\) having the
property specified in \hyperref[prop:II.2.g]{II.2(g)}.  If
\(n_1,n_2,\ldots\) is a sequence tending to infinity and \(f\) is an
arbitrary element of \(\mathfrak H\), then the sequence
\[
 F_{n_1}(x),F_{n_2}(x),\ldots
\]
has a subsequence \(F_{n_{p_1}}(x),F_{n_{p_2}}(x),\ldots\),
\(p_1<p_2<\cdots\), such that, for this \(f\), the set \(\mathfrak M\) of
all \(x\) on which it converges to \(F(x)\) has the property specified in
\hyperref[prop:II.2.g]{II.2(g)}.
\end{itshape}

% Source PDF page 21 / journal p. 211
\noindent\textit{Proof.}
The second half follows at once from \hyperref[thm:4]{Theorem~4} if the
sequence \(F_1(x),F_2(x),\ldots\) occurring there is replaced by
\[
 F_1(x),F(x),F_2(x),F(x),\ldots.
\]
The first half follows from the second if we can exhibit a function \(F(x)\)
having the properties named in the second half.  It should again be noted
that it is enough to restrict attention to a complete normalized
orthogonal system, that is, to a countable set of vectors
\(f_1,f_2,\ldots\).  This is just as permissible in the second half of the
present theorem as it was in \hyperref[thm:4.gamma]{\(\gamma\)} of
\hyperref[thm:4]{Theorem~4}, from which that half was obtained.  Our
starting point is the convergence of the \(F_n(A)\), and hence the
properties of the \(F_n(x)\) stated in
\hyperref[thm:4.gamma]{\(\gamma\)} of
\hyperref[thm:4]{Theorem~4}.

First let \(f\) be arbitrary but fixed, and let \(\epsilon>0\),
\(\eta>0\).  There must be a \(v_0=v_0(\epsilon,\eta)\) such that, for
every given pair \(m,n\geq v_0\), the set\tnref{C}
\[
 \mathfrak M_{m,n}(\epsilon)
 =\{x:|F_m(x)-F_n(x)|<\epsilon\}
\]
contains every \(\rho(\Psi_f(y))\), with the exception of a set of
\(y\)'s of Lebesgue measure at most \(\eta\).  Otherwise, there would be
two sequences \(m_1,m_2,\ldots\) and \(n_1,n_2,\ldots\), both tending to
infinity, such that, for every
\(\mathfrak M_{m_v,n_v}(\epsilon)\), \(v=1,2,\ldots\), the set of
\(y\)'s for which \(\rho(\Psi_f(y))\) does not belong to that set has
Lebesgue measure at least \(\eta\).  From the sequence of pairs
\((m_v,n_v)\), choose the subsequence \((m_{p_v},n_{p_v})\),
\(p_1<p_2<\cdots\), described in
\hyperref[thm:4.gamma]{\(\gamma\)} of
\hyperref[thm:4]{Theorem~4}.

All \(y\)'s under consideration lie in a
set of finite measure--namely \(0\leq y\leq|f|^2\), the range of values of
\(\Phi_f(x)=|E(\rho(x))f|^2\).  The Arzel\`a--Young theorem\noteref{39}
therefore applies: the set of \(y\)'s for which
\(\rho(\Psi_f(y))\notin\mathfrak M_{m_{p_v},n_{p_v}}(\epsilon)\) for
infinitely many \(v\) also has measure at least \(\eta\).  For such an
\(x=\rho(\Psi_f(y))\), the inequality
\[
 |F_{m_{p_v}}(x)-F_{n_{p_v}}(x)|\geq\epsilon
\]
holds infinitely often, so certainly
\(F_{m_{p_v}}(x)-F_{n_{p_v}}(x)\not\to0\).  Yet by
\hyperref[thm:4.gamma]{\(\gamma\)} of
\hyperref[thm:4]{Theorem~4}, the \(y\)'s whose
\(x=\rho(\Psi_f(y))\) produces a sequence
\(F_{m_{p_v}}(x)-F_{n_{p_v}}(x)\) not converging to \(0\) should form a
set of Lebesgue measure \(0\).  This is a contradiction.

Now let \(n_1,n_2,\ldots\) be a sequence tending to infinity.  Choose
integers \(p_1<p_2<p_3<\cdots\) so that
\begin{align*}
 n_{p_1}&\geq v_0\left(\frac12,\frac12\right),\\
 n_{p_2}&\geq
 \max\left\{v_0\left(\frac12,\frac12\right),
              v_0\left(\frac14,\frac14\right)\right\},\\
 n_{p_3}&\geq
 \max\left\{v_0\left(\frac14,\frac14\right),
              v_0\left(\frac18,\frac18\right)\right\},
 \qquad\ldots.
\end{align*}
Then
\[
 |F_{n_{p_v}}(x)-F_{n_{p_{v+1}}}(x)|<\frac1{2^v}
\]
for every \(\rho(\Psi_f(y))\), with the exception of a set of \(y\)'s of
Lebesgue measure at most \(2^{-v}\).  Thus, for all \(v\geq\mu\)
simultaneously (with \(\mu\) fixed), the preceding inequalities hold
except on a set of \(y\)'s of Lebesgue measure at most
\[
 \frac1{2^\mu}+\frac1{2^{\mu+1}}+\cdots=\frac1{2^{\mu-1}}.
\]
For every remaining \(y\), and for its
\(x=\rho(\Psi_f(y))\), if \(q>v\geq\mu\), then

% Source PDF page 22 / journal p. 212
\[
 |F_{n_{p_v}}(x)-F_{n_{p_q}}(x)|
 \leq\frac1{2^v}+\frac1{2^{v+1}}+\cdots+\frac1{2^{q-1}}
 <\frac1{2^{v-1}}.
\]
The \(F_{n_{p_v}}(x)\) therefore satisfy Cauchy's convergence criterion;
that is, the sequence \(F_{n_{p_v}}(x)\), \(v=1,2,\ldots\), converges.
The set of \(y\)'s for whose \(x=\rho(\Psi_f(y))\) it diverges thus has
measure at most \(2^{1-\mu}\), and, since this is true for every \(\mu\),
that measure is \(0\).

We now let \(f\) vary.  Take a sequence \(n_1,n_2,\ldots\) tending to
infinity and proceed as follows.  From
\(F_{n_1}(x),F_{n_2}(x),\ldots\), choose as above a subsequence
\(F_{r_1}(x),F_{r_2}(x),\ldots\) that converges for every
\(\rho(\Psi_{f_1}(y))\), except for a set of \(y\)'s of measure \(0\).
From \(F_{r_2}(x),F_{r_3}(x),\ldots\), choose a subsequence
\(F_{s_1}(x),F_{s_2}(x),\ldots\) that converges for every
\(\rho(\Psi_{f_2}(y))\), except for a set of \(y\)'s of measure \(0\).
From \(F_{s_2}(x),F_{s_3}(x),\ldots\), choose a subsequence
\(F_{t_1}(x),F_{t_2}(x),\ldots\) that converges for every
\(\rho(\Psi_{f_3}(y))\), except for a set of \(y\)'s of measure \(0\),
and continue in this fashion.

The diagonal sequence
\[
 F_{r_1}(x),F_{s_1}(x),F_{t_1}(x),\ldots
\]
is a subsequence of \(F_{n_1}(x),F_{n_2}(x),\ldots\), and, after finitely
many of its terms have been removed, it is a subsequence of each of the
sequences
\[
 \begin{gathered}
 F_{r_1}(x),F_{r_2}(x),\ldots,\qquad
 F_{s_1}(x),F_{s_2}(x),\ldots,\\
 F_{t_1}(x),F_{t_2}(x),\ldots,\quad\ldots.
 \end{gathered}
\]
Let \(\mathfrak M\) be the set of all \(x\)'s for which the diagonal
sequence converges.  By the last observation, \(\mathfrak M\) contains
every \(\rho(\Psi_{f_n}(y))\), except for a set of \(y\)'s of measure
\(0\), and this holds for every \(n\).  Instead of
\(F_{r_1}(x),F_{s_1}(x),\ldots\), write
\(F_{n_{p_1}}(x),F_{n_{p_2}}(x),\ldots\), \(p_1<p_2<\cdots\), and call
its limit function \(F(x)\).  For the moment, \(F\) is defined only on
\(\mathfrak M\); outside \(\mathfrak M\), set it equal, say, to \(0\).
Because the \(F_{n_{p_v}}(x)\) are uniformly bounded, so is \(F(x)\).
Because they are \(\Phi\)-measurable, so are \(\mathfrak M\) and \(F(x)\)
(this is seen exactly as for ordinary measurability).

Our construction still appears to make \(F(x)\) depend on the sequence
\(n_1,n_2,\ldots\).  Let \(m_1,m_2,\ldots\) and \(n_1,n_2,\ldots\) be
two such sequences, and let \(m_{p_1},m_{p_2},\ldots\) and
\(n_{q_1},n_{q_2},\ldots\), respectively, be the subsequences constructed
above, so that
\[
 F_{m_{p_v}}(x)\longrightarrow F'(x),
 \qquad
 F_{n_{q_v}}(x)\longrightarrow F''(x)
 \qquad(v\to\infty),
\]
both for every \(x=\rho(\Psi_{f_l}(y))\), except for a set of \(y\)'s of
measure \(0\), and for every \(l\).  By
\hyperref[thm:4.gamma]{\(\gamma\)} of
\hyperref[thm:4]{Theorem~4}, if \(l\) is fixed for the moment, there is a
sequence \(u_1<u_2<\cdots\), depending on \(l\), such that
\[
 F_{m_{p_{u_v}}}(x)-F_{n_{q_{u_v}}}(x)\longrightarrow0
\]
for every \(\rho(\Psi_{f_l}(y))\), with the exception of a set of \(y\)'s
of measure \(0\).  Hence \(F'(x)=F''(x)\) for every such
\(\rho(\Psi_{f_l}(y))\), apart from a null set of \(y\)'s, and this holds
for every \(l\).  Thus \(F(x)\) is essentially independent of the sequence
\(n_1,n_2,\ldots\): once chosen and fixed for one such sequence, it has
the properties set out in the second half of the theorem for every other
sequence as well.

% Source PDF page 23 / journal p. 213
This supplies everything that was still required to prove the first half
of the theorem.

Before closing this part, we mention that neither
\hyperref[thm:4]{Theorem~4} nor \hyperref[thm:5]{Theorem~5} remains true
for the ``weak'' topology of \(\mathbf B\) (cf.\ \noteref{35}); we omit
the simple counterexamples.  We further observe that a necessary and
sufficient criterion for \(F(A)=G(A)\) follows easily from the second half
of \hyperref[thm:5]{Theorem~5}: put all \(F_n(x)\) equal to \(G(x)\).
We obtain the following.  For every \(f\), the equality \(F(x)=G(x)\)
must hold at every \(\rho(\Psi_f(y))\), with the exception of a set of
\(y\)'s of Lebesgue measure \(0\).  Thus the sufficient condition in
\hyperref[prop:II.2.g]{II.2(g)} is also necessary.  As in
\hyperref[thm:5]{Theorem~5}, it is enough to require this only for a
complete normalized orthogonal system rather than for all \(f\).

Finally, the following point deserves emphasis.  In the second halves of
\hyperref[thm:4.gamma]{\(\gamma\)} in \hyperref[thm:4]{Theorem~4} and of
\hyperref[thm:5]{Theorem~5}, the subsequence occurring there,
\[
 F_{m_{p_v}}(x)-F_{n_{p_v}}(x)
 \quad\text{or}\quad
 F_{n_{p_v}}(x)-F(x),
 \qquad v=1,2,\ldots,
\]
is chosen in dependence on \(f\), and its set \(\mathfrak M\) has the
property of \hyperref[prop:II.2.g]{II.2(g)} only for this \(f\).  At both
places we may equally well require that the subsequence be chosen
independently of \(f\); its fixed set \(\mathfrak M\) then has the property
of \hyperref[prop:II.2.g]{II.2(g)} in general.  Indeed, if
\(f_1,f_2,\ldots\) is a complete normalized orthogonal system, a Hilbert
diagonal procedure, exactly as in the proof of
\hyperref[thm:5]{Theorem~5}, gives a subsequence (that is, a set
\(\mathfrak M\)) with the desired property for every
\(f\in\{f_1,f_2,\ldots\}\).  That such a set \(\mathfrak M\) has the
desired property for every \(f\) whatsoever follows, for example, from
the criterion for \(F(A)=G(A)\) just obtained; there too it was immaterial
whether the requirement on the (plainly general) set \(\mathfrak M\) was
made for all \(f\), or only for \(f_1,f_2,\ldots\).

\section*{III. Operator Functions and Rings}
\phantomsection\label{sec:III}

\subsection*{1.}\phantomsection\label{sec:III.1}

\hyperref[thm:4]{Theorems~4} and \hyperref[thm:5]{5} make possible a
simple proof of the assertion stated in \hyperref[sec:intro.3]{the
Introduction, 3}.

\medskip
\phantomsection\label{thm:6}\noindent\textbf{Theorem 6.}
\emph{Let \(A\) be a bounded Hermitian operator and \(\mathbf R(A)\) its
ring.  Then \(\mathbf R(A)\) is the set of all \(F(A)\), where \(F(x)\)
ranges over all \(\Phi\)-measurable functions satisfying \(F(0)=0\).
(It is immaterial whether the \(F(x)\) are required to be bounded, or, as
in \noteref{28}, merely bounded on every finite interval.)}

\noindent\textit{Proof.}
First, every such \(F(A)\) belongs to \(\mathbf R(A)\).  Indeed, as we have
repeatedly concluded, it is enough to consider functions \(F(x)\) of the
third Baire class, naturally with \(F(0)=0\).  By the closure properties
of rings (together with \hyperref[prop:II.2.h]{II.2(h)} or
\hyperref[thm:5]{Theorem~5}), it is then enough to consider even those of
the zeroth Baire class; here these may be taken to be the polynomials

% Source PDF page 24 / journal p. 214
\(p(x)\) with \(p(0)=0\).  We have long known that such a \(p(A)\) belongs
to the ring of \(A\), for example from \hyperref[ref:A]{A}, p.~404.

For the converse, recall that, by the observation at the same place,
every operator of \(\mathbf R(A)\) is the limit of a strongly doubly
convergent sequence
\[
 p_1(A),p_2(A),\ldots,
\]
where every \(p_n(x)\) is a polynomial and \(p_n(0)=0\).  It follows from
\hyperref[thm:5]{Theorem~5} that the operator is of the form \(F(A)\); the
characterization of \(F(x)\) given there shows that \(F\) is bounded and
\(\Phi\)-measurable and that it may be chosen with \(F(0)=0\).  This proves
the theorem.

Two remarks suggest themselves.  First, it is enough to restrict attention
to functions \(F(x)\) of the third Baire class.  Second, the restriction
\(F(0)=0\) may be dropped if altering \(F(x)\) at the point \(x=0\) does
not change \(F(A)\).  This means that, for every \(f\), the equality
\(\rho(\Psi_f(y))=0\) occurs only on a Lebesgue null set; in other words,
\(0\) is not the value of this function on an interval of constancy.
Equivalently, \(\rho^{-1}(0)\) is not a constancy-interval value of
\(\Psi_f(y)\), or \(\rho^{-1}(0)\) is not a discontinuity point of
\[
 \Phi_f(x)=|E(\rho(x))f|^2,
\]
or \(0\) is not a discontinuity point of \(|E(\xi)f|^2\).  Thus this is
the case precisely when \(0\) does not belong to the point spectrum
\(\mathfrak P\) of \(A\), or, equivalently, when
\[
 Af=0\quad\Longrightarrow\quad f=0.
\]

The Hermitian operators in the ring \(\mathbf R(A)\) are obtained by
restricting \(F(x)\) to real-valued functions.  It is clear that \(F(A)\)
is Hermitian when \(F(x)\) is real.  Conversely, if \(F(A)\) is Hermitian,
then \(F(A)=F(A)^*\), and hence
\[
 F(A)=\frac{F(A)+F(A)^*}{2}=G(A),
\]
where \(G(x)\) is the real function
\[
 G(x)=\frac{F(x)+\overline{F(x)}}{2}.
\]

An Abelian set \(\mathbf M\) of operators in \(\mathbf B\)--that is, a
set such that, for any two of its elements \(A,B\), the operator \(A\)
commutes under multiplication with both \(B\) and \(B^*\), and whose
elements are therefore all normal--is always contained in an Abelian ring
(cf.\ \hyperref[ref:A]{A}, p.~389).  By the principal result of
\hyperref[ref:A]{A} (p.~401, Theorem~10), it is consequently contained in
a ring \(\mathbf R(A)\).  Thus there certainly exists a Hermitian operator
\(A\) of which all elements of \(\mathbf M\) are functions.  If
\(\mathbf M\) consists of Hermitian operators, it is Abelian whenever
they commute pairwise under multiplication, and the functions just
mentioned may then even be chosen real.

One further point is worth mentioning.  Let \(A\) be a Hermitian operator,
and consider those operators \(B\in \mathbf B\) that commute with every
operator commuting with \(A\).  It was proved in \hyperref[ref:A]{A},
p.~404, that all these \(B\) are normal; indeed, they are precisely the
operators
\[
 a\,1+B',
\]
where \(a\) ranges over all complex numbers and \(B'\) over all elements

% Source PDF page 25 / journal p. 215
of $\mathbf R(A)$.  By \hyperref[thm:6]{Theorem~6}, they are therefore precisely all
\[
  a\,1+F'(A)=F(A), \qquad F(x)=a+F'(x),
\]
where $a$ ranges over all complex numbers and $F'(x)$ over all functions satisfying
$F'(0)=0$---that is, $F(x)$ ranges over all functions.  If we assume $B$ to be
Hermitian, it again suffices to consider real-valued $F(x)$.

\subsection*{2.}
\phantomsection\label{sec:III.2}
Let $\mathbf R$ be an Abelian ring.  It can then always be written as
$\mathbf R(A)$ (cf. A, p.~401, Theorem~10), with $A$ Hermitian; as is easily
seen, however, this by no means determines $A$.  Suppose, then, that
\[
  \mathbf R=\mathbf R(A)=\mathbf R(B).
\]
What relation must hold between $A$ and $B$ (both assumed Hermitian)?

Since each belongs to the ring of the other, each is a function---indeed, a real
function---of the other:
\[
  B=F(A), \qquad A=G(B), \qquad F(0)=G(0)=0.
\]
Conversely, this is plainly also sufficient for $A$ to belong to the ring of $B$ and
$B$ to the ring of $A$, that is, for $\mathbf R(A)=\mathbf R(B)$.  We now
regard the function $F(x)$ (which is, of course, bounded in every finite interval and
$\Phi$-measurable), with $F(0)=0$, as given, and ask: what conditions must it satisfy
in order that there exist a function $G(x)$ (likewise bounded in every finite interval
and $\Phi$-measurable), with $G(0)=0$, such that, for the operator $B$ defined by
$B=F(A)$, one has $A=G(B)$?

This requirement says that $G(F(A))=A$, hence, by the concluding remark of
\hyperref[sec:II.3]{II, 3}, that $G(F(x))=x$ on an $x$-set $\mathfrak{M}$ having the
property specified in \hyperref[sec:II.2]{II, 2, g)}.  As one sees, this amounts to
requiring the $\Phi$-measurable function $F(x)$ to possess a $\Phi$-measurable inverse
$G(x)$, except on an $x$-set whose inverse images under every mapping
$\varrho(\psi_f(y))=x$ are null sets (cf. \hyperref[sec:II.2]{II, 2, g)}).  (The
boundedness of $G(x)$ is immaterial: only the $x$ in the spectrum of $A$ matter, hence
only those in a finite interval; consequently, wherever the range of values of $G(u)$
is prescribed at all, it lies between finite bounds.)  It would be natural to try to
dispense with the assumption, and thus, for example, to work with functions $F(x)$
that assume every real value once and only once.  But then one would still have to
settle the further question: when is the inverse of a $\Phi$-measurable function (that
has an inverse) itself $\Phi$-measurable?  This question appears to be rather
complicated.\noteref{40}

\section*{IV. Functions of Several Variables and Normal Operators}
\phantomsection\label{sec:IV}

\subsection*{1.}
\phantomsection\label{sec:IV.1}
We first introduce the notion of $\Phi$-measurability for (complex-valued) functions
$F(x_1,x_2,\ldots)$ that depend on a finite or countably infinite number of complex
variables $x_1,x_2,\ldots$ (even when there are finitely many variables, in which case
they should be denoted $x_1,\ldots,x_n$,

% Source PDF page 26 / journal p. 216
we write $x_1,x_2,\ldots$ for uniformity).

We call such a function $F(x_1,x_2,\ldots)$ $\Phi$-measurable if it has the following
property: whenever $f_1(u),f_2(u),\ldots$ is an arbitrary system of functions defined
for every real $u$ and right-continuous,\noteref{41} then
$F(f_1(u),f_2(u),\ldots)$, regarded as a complex-valued function of the real variable
$u$, is measurable.  Clearly, every function that actually depends on only finitely
many variables and is continuous in them is $\Phi$-measurable.  Moreover, if
$F_1,F_2,\ldots$ is a sequence of $\Phi$-measurable functions converging everywhere
to $F$, then $F$ is also $\Phi$-measurable.  Rather than prove only this last assertion,
we prove at once the stronger one.  Let $F_1,F_2,\ldots$ be $\Phi$-measurable and
define
\[
 F(x_1,x_2,\ldots)=
 \begin{cases}
   \displaystyle\lim_{n\to\infty}F_n(x_1,x_2,\ldots),
      & \text{if this limit exists},\\[2pt]
   0, & \text{otherwise}.
 \end{cases}
\]
Indeed, since the $\Phi$-measurability of $F_n$ and $F$ means the ordinary
measurability of
\[
 F_n(f_1(u),f_2(u),\ldots), \qquad F(f_1(u),f_2(u),\ldots)
\]
for every choice of $f_1(u),f_2(u),\ldots$, it is enough to prove the assertion above
for functions of one real variable $u$ and ordinary measurability; there, however, it
is well known to be true.  Thus all functions of all Baire classes are again
$\Phi$-measurable.\noteref{42}  From the analogous properties of measurability for
functions of one real variable, it follows immediately that, with $F$ and $G$, the
functions $F\pm G$, $aF$, and $FG$ are also $\Phi$-measurable.  The following theorem
goes somewhat deeper.

\medskip
\phantomsection\label{thm:7}
\noindent\textbf{Theorem 7.}
\textit{If $F(x_1,x_2,\ldots)$ is $\Phi$-measurable and
$G_1(y_1,y_2,\ldots),G_2(y_1,y_2,\ldots),\ldots$ are measurable, respectively
$\Phi$-measurable (all the functions being complex-valued and having complex
arguments), then}
\[
 F\bigl(G_1(y_1,y_2,\ldots),G_2(y_1,y_2,\ldots),\ldots\bigr)
\]
\textit{is measurable, respectively $\Phi$-measurable.  (Observe that the first
assertion, which concerns ordinary measurability, has meaning only for finitely many
variables, since ordinary measurability, unlike $\Phi$-measurability, is defined only
in that case.)}

\smallskip
\noindent\textit{Proof.}
The second assertion ($\Phi$-measurability) amounts to saying that, for all functions
$f_1(u),f_2(u),\ldots$ right-continuous, the function
\[
 F\bigl(G_1(f_1(u),f_2(u),\ldots),
         G_2(f_1(u),f_2(u),\ldots),\ldots\bigr)
\]
is measurable.  But since
\[
 G_1(f_1(u),f_2(u),\ldots),
 \quad G_2(f_1(u),f_2(u),\ldots),\quad\ldots
\]
are measurable, this is precisely the first assertion for one real variable.  The
first assertion for finitely many complex variables $y_1,\ldots,y_n$ (its general
form) amounts to the same thing as for twice as many real variables
\[
  \Re y_1,\Im y_1,\ldots,\Re y_n,\Im y_n;
\]
hence only the first assertion for finitely many real variables remains to be proved.
We again denote these real variables by $y_1,y_2,\ldots$.  Next, by replacing the
variables $x_1,x_2,\ldots$ of $F$ by
\[
  \Re x_1,\Im x_1,\Re x_2,\Im x_2,\ldots
\]
and correspondingly replacing

% Source PDF page 27 / journal p. 217
$G_1,G_2,\ldots$ by
\[
 \Re(G_1),\Im(G_1),\Re(G_2),\Im(G_2),\ldots,
\]
we arrange that $F$ too has real arguments; we again call these arguments
$x_1,x_2,\ldots$, and again call the functions to be substituted
$G_1,G_2,\ldots$.  Finally, by the same device used at the beginning of the proof of
\hyperref[thm:2]{Theorem~2}, we may restrict all the variables $x_1,x_2,\ldots$ and
$y_1,\ldots,y_n$ to finite intervals.  In place of the interval $0<\xi<2$ used there,
we now choose $0<\xi<1$, which may of course be extended to $0\leq\xi<1$.  Thus
$F(x_1,x_2,\ldots)$ is defined for
\[
 0\leq x_1<1,\quad 0\leq x_2<1,\quad\ldots;
\]
the functions $G_1(y_1,\ldots,y_n),G_2(y_1,\ldots,y_n),\ldots$ are defined for
\[
 0\leq y_1<1,\quad\ldots,\quad 0\leq y_n<1,
\]
and the values of $G_1,G_2,\ldots$ (since they are substituted into $F$) always satisfy
$0\leq G_i<1$.

Suppose now that the proof has been carried out for $n=1$.  The general case with
$y_1,\ldots,y_n$ is reduced to the case of a single $y$ as follows.  Let
\[
 y=\alpha(y_1,\ldots,y_n),
 \qquad y_1=\beta_1(y),\ldots,y_n=\beta_n(y)
\]
be a one-to-one mapping of the interval $0\leq y<1$ onto the cube
\[
 0\leq y_1<1,\quad\ldots,\quad 0\leq y_n<1
\]
(a one-dimensional, respectively $n$-dimensional, Lebesgue null set may be omitted
from either side) that preserves measure, in the following sense: a $y$-preimage and
its $\{y_1,\ldots,y_n\}$-image are either both measurable or both nonmeasurable, and
in the former case have the same measure (in the one-dimensional and
$n$-dimensional senses, respectively).\noteref{43}  Then, together with the measurable
functions $G_1(y_1,\ldots,y_n),G_2(y_1,\ldots,y_n),\ldots$, the functions
\[
 G_1(\beta_1(y),\ldots,\beta_n(y)),
 \quad G_2(\beta_1(y),\ldots,\beta_n(y)),\quad\ldots
\]
are measurable.  Consequently, once the theorem is established for one variable,
\[
 F\bigl(G_1(\beta_1(y),\ldots,\beta_n(y)),
         G_2(\beta_1(y),\ldots,\beta_n(y)),\ldots\bigr)
\]
is measurable as well.  It remains measurable when $y$ in it is replaced by
$\alpha(y_1,\ldots,y_n)$; that is,\tnref{A}
\[
 F\bigl(G_1(y_1,\ldots,y_n),G_2(y_1,\ldots,y_n),\ldots\bigr)
\]
is measurable.  We therefore pass to the case of a single real variable $y$.

Let $f_1(u),f_2(u),\ldots$ be right-continuous, and let $\varphi(y)$ be
monotone nondecreasing and right-continuous.  Then all of
$f_1(\varphi(y)),f_2(\varphi(y)),\ldots$ are right-continuous, so
\[
 F\bigl(f_1(\varphi(y)),f_2(\varphi(y)),\ldots\bigr)
\]
is measurable.  Since this holds for every such $\varphi(y)$,
$F(f_1(u),f_2(u),\ldots)$ is $\Phi$-measurable; hence, by
\hyperref[thm:2]{Theorem~2}, for every measurable $G(y)$ the function
\[
 F\bigl(f_1(G(y)),f_2(G(y)),\ldots\bigr)
\]
is measurable.  If, for every system of measurable functions
$G_1(y),G_2(y),\ldots$, one can now find a measurable function $G(y)$ and functions
$f_1(u),f_2(u),\ldots$ right-continuous such that
\[
 G_1(y)=f_1(G(y)),\qquad G_2(y)=f_2(G(y)),\quad\ldots,
\]
then the proof is complete.

We carry out this construction as follows.  Let
\[
 x=\alpha(x_1,x_2,\ldots),
 \qquad x_1=\beta_1(x),\quad x_2=\beta_2(x),\quad\ldots
\]
be a one-to-one mapping of a subset of the interval $0\leq x<1$ onto the entire cube
$0\leq x_1<1$, $0\leq x_2<1$, \ldots.  Let all the functions
$\beta_1(x),\beta_2(x),\ldots$ be right-continuous, and let them have the
following property: when $x$ runs through the part of the image set lying in an
interval
\[
 \frac{p}{2^q}\leq x<\frac{p+1}{2^q}
 \qquad(q=0,1,\ldots;\;p=0,1,\ldots,2^q-1),
\]
the point $x_1=\beta_1(x),x_2=\beta_2(x),\ldots$ runs through

% Source PDF page 28 / journal p. 218
a cube (that is, a set $a_1\leq x_1<b_1$, $a_2\leq x_2<b_2$, \ldots).\noteref{44}
Now put
\[
 G(y)=\alpha(G_1(y),G_2(y),\ldots).
\]
The condition
\[
 \frac{p}{2^q}\leq G(y)<\frac{p+1}{2^q}
\]
is equivalent to
\[
 a_1\leq G_1(y)<b_1,\qquad a_2\leq G_2(y)<b_2,\qquad\ldots;
\]
hence it holds on a $y$-set that is the intersection of countably many measurable
sets (since $G_1(y),G_2(y),\ldots$ are measurable), and is therefore itself
measurable.  The $y$-set on which $G(y)<a$, for arbitrary $a$, can be formed as a
union of countably many such sets, and is therefore measurable as well.  Thus $G(y)$
is measurable.  Taking $f_1=\beta_1,f_2=\beta_2,\ldots$, we have reached our goal.

It follows at once from \hyperref[thm:7]{Theorem~7} that, if the number of variables
$x_1,x_2,\ldots$ is finite, then every $\Phi$-measurable
$F(x_1,x_2,\ldots)$ is measurable: it is enough to identify
$y_1,y_2,\ldots$ with $x_1,x_2,\ldots$ and to put
\[
 G_1(y_1,y_2,\ldots)=y_1,\qquad
 G_2(y_1,y_2,\ldots)=y_2,\quad\ldots.
\]
Even for a single variable $x$, the new notion of $\Phi$-measurability cannot be
compared directly with the old one (\hyperref[sec:I.3]{I, 3}), because $x$ is now
complex whereas before it was real.  If, however, $F(x)$ depends only on $\Re x$, the
two notions are equivalent: by their definitions, the old $\Phi$-measurability
follows from the new, while the converse follows from \hyperref[thm:2]{Theorem~2} and
the definition of the new $\Phi$-measurability.

If $F(x_1,x_2,\ldots)$ is $\Phi$-measurable in the new sense and
$G_1(x),G_2(x),\ldots$ are $\Phi$-measurable in the old sense (with $x$ real), then,
by what has just been said, $G_1(\Re x),G_2(\Re x),\ldots$ (with $x$ complex) are
$\Phi$-measurable in the new sense.  By \hyperref[thm:7]{Theorem~7}, therefore,
\[
 F\bigl(G_1(\Re x),G_2(\Re x),\ldots\bigr)
\]
is $\Phi$-measurable in the new sense; for real $x$ this is again
$F(G_1(x),G_2(x),\ldots)$ in the old sense.

In view of what has been said, we may henceforth drop the distinction between the two
definitions of $\Phi$-measurability without danger of confusion.

\subsection*{2.}
\phantomsection\label{sec:IV.2}
Let $A_1,A_2,\ldots$ now be a finite or countably infinite Abelian set of bounded
linear operators (cf. A, p.~389); that is, every $A_m$ is to commute with every $A_n$
and with $A_n^*$---in particular, all the $A_n$ are normal.  Let
$F(x_1,x_2,\ldots)$ be a $\Phi$-measurable function with as many variables as there
are operators $A_1,A_2,\ldots$.

According to one of the remarks in \hyperref[sec:III.1]{III, 1}, there exist bounded
Hermitian operators of which all $A_1,A_2,\ldots$ are functions.  Let $A$ be one such
operator, so that
\[
 A_1=G_1(A),\qquad A_2=G_2(A),\quad\ldots,
\]
where $G_1(x),G_2(x),\ldots$ are all $\Phi$-measurable (and $x$ is real).  The
function
\[
 H(x)=F(G_1(x),G_2(x),\ldots)
\]
is consequently also $\Phi$-measurable.  If it is bounded on a set $\mathfrak{M}$
having the properties specified in \hyperref[sec:II.2]{II, 2, g)} (for the operator
$A$), then we can form $H(A)$; this is a bounded linear normal operator.  We shall
define it to be $F(A_1,A_2,\ldots)$.  For this purpose, however, we must first show
that the dependence of this definition on the auxiliary operator $A$ is only apparent:
that is, that the definition is possible for all the indicated choices of $A$, or

% Source PDF page 29 / journal p. 219
for none of them, and that in the former case it gives the same result for all of them.

The ring $\mathbf R(A_1,A_2,\ldots)$ is Abelian and hence, by A, p.~401,
Theorem~10, is an $\mathbf R(B)$ with $B$ Hermitian.  By
\hyperref[thm:6]{Theorem~6}, all $A_1,A_2,\ldots$ are therefore functions of $B$;
thus $B$ may be used as an $A$ in the preceding sense.  We shall now show that every
one of the indicated operators $A$ behaves, for the purposes of our definition,
exactly as this $B$ does and leads to the same result.  This will establish everything
required.

We have
\[
 G_1(A)=A_1,\quad G_2(A)=A_2,\quad\ldots,
 \qquad
 G'_1(B)=A_1,\quad G'_2(B)=A_2,\quad\ldots.
\]
The ring $\mathbf R(A)$ contains all functions of $A$, hence also
$A_1,A_2,\ldots$; it therefore contains
$\mathbf R(A_1,A_2,\ldots)=\mathbf R(B)$, and so contains $B$.  Consequently
$B=K(A)$.  It follows that
\[
 G_1(A)=G'_1(K(A)),\qquad G_2(A)=G'_2(K(A)),\quad\ldots.
\]
By \hyperref[sec:II.2]{II, 2, f)} and the concluding remark of
\hyperref[sec:II.3]{II, 3}, therefore,
\[
 G_1(x)=G'_1(K(x))
\]
on a set $\mathfrak{M}$ having the properties of \hyperref[sec:II.2]{II, 2, g)} for
$A$, and likewise $G_2(x)=G'_2(K(x))$ on such a set, and so on.  All these equalities
hold simultaneously on the intersection of this sequence of sets, which is again a
set $\mathfrak{M}$ of the kind specified in \hyperref[sec:II.2]{II, 2, g)}.  Under our
definition we now have, for $A$, to form
\[
 H(x)=F(G_1(x),G_2(x),\ldots),
\]
and, for $B$, to form
\[
 H'(x)=F(G'_1(x),G'_2(x),\ldots).
\]
By the preceding result, $H(x)=H'(K(x))$ on a set $\mathfrak{M}$ of the kind
specified in \hyperref[sec:II.2]{II, 2, g)}.  It follows at once from
\hyperref[sec:II.2]{II, 2, f)} that
\[
 H'(B)=H'(K(A))=H(A).
\]
Thus the two routes to the definition give the same result, provided that $H(x)$ for
$A$ and $H'(x)$ for $B$ have the desired boundedness properties.  It remains to prove
that these properties are equivalent.  We shall in fact prove the stronger assertion
that the validity of $|H(x)|\leq C$ on a set $\mathfrak{M}$ of the kind specified in
\hyperref[sec:II.2]{II, 2, g)} for $A$ is equivalent to that of
$|H'(x)|\leq C$ on such a set for $B$.

Replace $F(x_1,x_2,\ldots)$ by
\[
 F_D(x_1,x_2,\ldots)=
 \begin{cases}
   F(x_1,x_2,\ldots), & |F(x_1,x_2,\ldots)|\leq D,\\
   D, & \text{otherwise}.
 \end{cases}
\]
Then $H(x)$ and $H'(x)$ become, respectively,\tnref{B}
\[
 H_D(x)=
 \begin{cases}
   H(x), & |H(x)|\leq D,\\
   D, & \text{otherwise},
 \end{cases}
 \qquad
 H'_D(x)=
 \begin{cases}
   H'(x), & |H'(x)|\leq D,\\
   D, & \text{otherwise}.
 \end{cases}
\]
Since $H_D$ and $H'_D$ are always bounded, $H_D(A)$ and $H'_D(B)$ can certainly be
formed, and they are equal.  The stated boundedness property of $H(x)$, respectively
$H'(x)$, now says that
\[
 H_C(x)=H_{C+1}(x)
\]
on a set $\mathfrak{M}$ of the kind specified in \hyperref[sec:II.2]{II, 2, g)} for
$A$, respectively that $H'_C(x)=H'_{C+1}(x)$

% Source PDF page 30 / journal p. 220
on a set $\mathfrak{M}$ of that kind for $B$.  By the concluding remark of
\hyperref[sec:II.3]{II, 3}, this is equivalent to
\[
 H_C(A)=H_{C+1}(A),
 \qquad\text{respectively}\qquad
 H'_C(B)=H'_{C+1}(B).
\]
But the two left-hand sides are equal to one another, and so are the two right-hand
sides; hence the two conditions are equivalent, as asserted.

This proves that our definition of $F(A_1,A_2,\ldots)$ is unambiguous.  It remains to
mention that, wherever the old definition applies, the new one is equivalent to it.
If $F(x)$ has one variable $x=x_1$, into which a Hermitian $A_1$ is to be substituted,
we may put $A=A_1$ and $G_1(x)=x$ (so $G_1(A_1)=A_1=A$).  Then $F(A_1)$ in the new
sense is $F(A)=F(A_1)$ in the old sense, and the boundedness requirements (that is,
the conditions under which the expressions are meaningful) are identical as well.

\subsection*{3.}
\phantomsection\label{sec:IV.3}
We verify the fundamental properties of the new function concept.
\begin{enumerate}[label=\alph*)]
\item If
\[
 F(x_1,x_2,\ldots)=\overline{G(\overline{x}_1,\overline{x}_2,\ldots)},
\]
then
\[
 F(A_1,A_2,\ldots)=G(A_1^*,A_2^*,\ldots)^*.
\]

\item If $F(x_1,x_2,\ldots)=aG(x_1,x_2,\ldots)$, then
\[
 F(A_1,A_2,\ldots)=aG(A_1,A_2,\ldots).
\]

\item If
$F(x_1,x_2,\ldots)=G(x_1,x_2,\ldots)+H(x_1,x_2,\ldots)$, then
\[
 F(A_1,A_2,\ldots)=G(A_1,A_2,\ldots)+H(A_1,A_2,\ldots).
\]

\item If
$F(x_1,x_2,\ldots)=G(x_1,x_2,\ldots)H(x_1,x_2,\ldots)$, then
\[
 F(A_1,A_2,\ldots)=G(A_1,A_2,\ldots)H(A_1,A_2,\ldots).
\]
\end{enumerate}
(In a)--d), meaningfulness is to be understood as follows: whenever the right-hand
side is meaningful, the left-hand side is meaningful as well.)

\begin{enumerate}[label=\alph*),start=5]
\item If
\[
 F(x_1,x_2,\ldots)
 =G\bigl(H_1(x_1,x_2,\ldots),H_2(x_1,x_2,\ldots),\ldots\bigr),
\]
then
\[
 F(A_1,A_2,\ldots)
 =G\bigl(H_1(A_1,A_2,\ldots),H_2(A_1,A_2,\ldots),\ldots\bigr).
\]
\end{enumerate}
Indeed, by a) and d), all the operators
$H_1(A_1,A_2,\ldots),H_2(A_1,A_2,\ldots),\ldots$, whenever they are meaningful at
all, are normal; in this case either both sides are meaningful or neither is.  All
these properties follow without difficulty by applying the definition of
$F(x_1,x_2,\ldots)$ and recalling the corresponding properties of the old function
concept (cf. \hyperref[sec:II.2]{II, 2, b)--f)}).

It is unnecessary to give precise criteria for equality and convergence under the new
function concept, like those we had under the old one
(\hyperref[sec:II.2]{II, 2, g) and h)}), or in greater detail in the concluding remark
of \hyperref[sec:II.3]{II, 3} and in \hyperref[thm:4]{Theorems~4} and
\hyperref[thm:5]{5}; in each case it suffices to reduce the new function concept to
the old by its definition.  Nevertheless, we record the following, crudest form of
the convergence theorem.

\begin{enumerate}[label=\alph*),start=6]
\item Suppose that everywhere
\[
 F_n(x_1,x_2,\ldots)\longrightarrow F(x_1,x_2,\ldots),
\]
and that the $F_n(x_1,x_2,\ldots)$ are uniformly bounded.  (Then
$F(x_1,x_2,\ldots)$ is also bounded, so all of
$F_n(A_1,A_2,\ldots)$ and $F(A_1,A_2,\ldots)$ are meaningful.)  Then, in the sense
of strong double convergence in $\mathbf B$,
\[
 F_n(A_1,A_2,\ldots)\longrightarrow F(A_1,A_2,\ldots).
\]
\end{enumerate}
This is evident, since the old function concept had this property as well.

\subsection*{4.}
\phantomsection\label{sec:IV.4}
We now prove the following strengthening of the penultimate assertion of
\hyperref[sec:III.1]{III, 1}.

% Source PDF page 31 / journal p. 221
\medskip
\phantomsection\label{thm:8}
\noindent\textbf{Theorem 8.}
\textit{Let $\mathbf M$ be an Abelian set of operators in $\mathbf B$ (that is,
its elements are bounded linear operators and, for any two of them, $A,B$, the
operator $A$ commutes with $B$ and with $B^*$; in particular, all of them are normal).
Then there exists an operator $R$ such that, on the one hand, every $A\in \mathbf M$
is a function of $R$,}
\[
 A=F_A(R),
\]
\textit{and, on the other hand, there exists a finite or countably infinite subset
$A_1,A_2,\ldots$ of $\mathbf M$ such that $R$ is a function of
$A_1,A_2,\ldots$,}
\[
 R=G(A_1,A_2,\ldots).
\]
\textit{Moreover, $R$ may be chosen Hermitian, and all the functions may be chosen to
vanish at zero:}
\[
 F_A(0)=0\quad(A\in \mathbf M),
 \qquad G(0,0,\ldots)=0.
\]

\smallskip
\noindent\textit{Proof.}
The ring $\mathbf R(\mathbf M)$, like $\mathbf M$ itself, is Abelian and may
therefore (as we have mentioned several times) be written as $\mathbf R(R)$, with
$R$ Hermitian.  Since every $A\in \mathbf M$ thus belongs to
$\mathbf R(\mathbf M)=\mathbf R(R)$, it is, by
\hyperref[thm:6]{Theorem~6}, equal to $F_A(R)$, where $F_A(0)=0$.  Conversely,
$R$ belongs to $\mathbf R(R)=\mathbf R(\mathbf M)$ and hence, by A, p.~398,
Theorem~9, is the limit of a strongly doubly convergent sequence of expressions, each
of which is obtained from elements of $\mathbf M$ by finitely many operations
\[
 A^*,\qquad aA,\qquad A+B,\qquad AB.
\]
Only finitely many elements of $\mathbf M$ can enter each of these expressions, so
only finitely or countably infinitely many occur in the whole sequence.  Denote them
by $A_1,A_2,\ldots$.  Each term of the sequence is therefore a polynomial, with no
constant term, in finitely many of
\[
 A_1,A_2,\ldots,A_1^*,A_2^*,\ldots.
\]
If we replace $A_1,A_2,\ldots$ by the variables $x_1,x_2,\ldots$ and
$A_1^*,A_2^*,\ldots$ by their complex conjugates
$\overline{x}_1,\overline{x}_2,\ldots$, call the $n$th such polynomial
$p_n(x_1,x_2,\ldots)$.  By Note~\noteref{42}, $p_n(x_1,x_2,\ldots)$ is
$\Phi$-measurable; and, since there is no constant term,
$p_n(0,0,\ldots)=0$.  By \hyperref[sec:II.2]{II, 2, a)} and
\hyperref[sec:IV.2]{IV, 2, a)--d)}, the $n$th term of our operator sequence is then
$p_n(A_1,A_2,\ldots)$, and consequently
\[
 p_n(A_1,A_2,\ldots)\longrightarrow R.
\]
It follows that
\[
 p_n\bigl(F_{A_1}(R),F_{A_2}(R),\ldots\bigr)\longrightarrow R.
\]
By \hyperref[thm:5]{Theorem~5}, therefore, there is a subsequence
$n_1<n_2<\cdots$ of $1,2,\ldots$ such that
\[
 p_{n_\nu}\bigl(F_{A_1}(x),F_{A_2}(x),\ldots\bigr)\longrightarrow x
 \qquad(\nu\to\infty)
\]
on an $x$-set $\mathfrak M$ of the kind specified in
\hyperref[sec:II.2]{II, 2, g)} for $R$.

Now put
\[
 G(x_1,x_2,\ldots)=
 \begin{cases}
   \displaystyle\lim_{\nu\to\infty}p_{n_\nu}(x_1,x_2,\ldots),
     & \text{if this limit exists},\\[2pt]
   0, & \text{otherwise}.
 \end{cases}
\]
By Note~\noteref{42}, this $G(x_1,x_2,\ldots)$ is $\Phi$-measurable, and by what has
just been observed,
\[
 G\bigl(F_{A_1}(x),F_{A_2}(x),\ldots\bigr)=x
\]
on a set $\mathfrak M$ of the kind specified in
\hyperref[sec:II.2]{II, 2, g)} for $R$.  Hence, by definition,
\[
 G(A_1,A_2,\ldots)=R.
\]
Since $p_{n_\nu}(0,0,\ldots)=0$, it follows that $G(0,0,\ldots)=0$.

Furthermore:

\medskip
\phantomsection\label{thm:9}
\noindent\textbf{Theorem 9.}
\textit{Let $\mathbf M$ be as in \hyperref[thm:8]{Theorem~8}.  Then
$\mathbf R(\mathbf M)$ is the set of all $F(A_1,A_2,\ldots)$, where
$F(x_1,x_2,\ldots)$ ranges over all $\Phi$-measurable functions satisfying}
\[
 F(0,0,\ldots)=0,
\]
\textit{and $A_1,A_2,\ldots$ may be any finite or countably infinite%
% Source PDF page 32 / journal p. 222
subset of $\mathbf M$.  The result remains true if only one such set
$A_1,A_2,\ldots$ is admitted, provided it is chosen suitably (as is always possible).}

\smallskip
\noindent\textit{Proof.}
Let $R,F_A(x),A_1,A_2,\ldots,G(x_1,x_2,\ldots)$ be as in the proof of
\hyperref[thm:8]{Theorem~8}.  If $B_1,B_2,\ldots$ are any finitely or countably
infinitely many elements of $\mathbf M$, then, since
\[
 B_1=F_{B_1}(R),\qquad B_2=F_{B_2}(R),\quad\ldots,
\]
every $F(B_1,B_2,\ldots)$ equals
\[
 F\bigl(F_{B_1}(R),F_{B_2}(R),\ldots\bigr)=H(R),
\]
where
\[
 H(x)=F\bigl(F_{B_1}(x),F_{B_2}(x),\ldots\bigr),
 \qquad H(0)=F(0,0,\ldots)=0.
\]
Thus it belongs to $\mathbf R(R)=\mathbf R(\mathbf M)$.  Conversely, every
element of $\mathbf R(\mathbf M)$ is an element of $\mathbf R(R)$, and is
therefore equal to $J(R)$ for some $J$ with $J(0)=0$.  It is consequently equal to
\[
 J\bigl(G(A_1,A_2,\ldots)\bigr)=K(A_1,A_2,\ldots),
\]
where
\[
 K(x_1,x_2,\ldots)=J\bigl(G(x_1,x_2,\ldots)\bigr),
 \qquad K(0)=J(0)=0.
\]
It therefore has the desired form, even for this one special choice of
$B_1,B_2,\ldots$, namely $B_1=A_1,B_2=A_2,\ldots$.

As is clear from the proof of \hyperref[thm:8]{Theorem~8}, the functions
$G(x_1,x_2,\ldots)$, respectively $F(x_1,x_2,\ldots)$, occurring in
\hyperref[thm:8]{Theorems~8} and \hyperref[thm:9]{9} may be restricted to the class
of those obtained, by a single passage to a limit in the sense of Note~\noteref{42},
from polynomials $p(x_1,x_2,\ldots)$ in finitely many variables and their complex
conjugates.  In the sense of Note~\noteref{42}, therefore, they are of the first Baire
class, even if the zeroth class is restricted to the stated polynomials
$p(x_1,x_2,\ldots)$.  If instead the Baire classes are defined by sequences of
functions converging everywhere (cf. \hyperref[sec:IV.2]{IV, 2, f)}), three
successive limiting processes are required to obtain these functions, as is easily
seen.  In this customary sense, therefore, all functions up to the third Baire class
are required.  Since every function may thus be replaced by one of this kind, the
same observation applies to the $F_A(x)$ in \hyperref[thm:8]{Theorem~8}.

Finally, we make an observation corresponding to the concluding remark of
\hyperref[sec:III.1]{III, 1}.  If the condition $F(0,0,\ldots)=0$ is omitted from the
statement of \hyperref[thm:9]{Theorem~9}, then in general we obtain not
$\mathbf R(\mathbf M)$, but the totality of all
\[
 a\,1+B,
\]
where $a$ is a complex number and $B\in \mathbf R(\mathbf M)$.  By A, p.~397,
Theorem~7, this is the set $\mathbf M''$, namely the set of all $A$ that commute
with every $C$ which commutes with $B$ and $B^*$ for every $B\in \mathbf M$.  Thus
these $A$ are precisely all the functions $F(A_1,A_2,\ldots)$ (with
$A_1,A_2,\ldots$ as in \hyperref[thm:9]{Theorem~9}).

\section*{Notes}
\phantomsection\label{sec:notes}

% Source PDF page 33 / journal p. 223
\originalnote{1}Throughout this paper two earlier works of the author will be cited
repeatedly: \phantomsection\label{ref:E}\emph{Allgemeine Eigenwerttheorie
Hermitescher Funktionaloperatoren} [``General Eigenvalue Theory of Hermitian
Functional Operators''], \emph{Mathematische Annalen}, vol.~102, no.~1,
pp.~49--131 (1929), and \phantomsection\label{ref:A}\emph{Zur Algebra der
Funktionaloperatoren und der Theorie der normalen Operatoren} [``On the Algebra
of Functional Operators and the Theory of Normal Operators''],
\emph{Mathematische Annalen}, vol.~102, no.~3, pp.~370--427 (1929).  We denote these
works by E and A, respectively.  The reference here is to A, in particular to the end
of Chapter~III, \S2.

It should also be pointed out that functions of operators had already been considered
by F.~Riesz and played, among other things, an important role in his proof of Hilbert's
theorem on the spectral form of bounded Hermitian operators.  More recently, M.~Stone
announced a general definition of functions of operators (\emph{Proceedings of the
National Academy of Sciences}, vol.~16, no.~2, pp.~172--175, 1930); moreover, special
constructions of operator functions play a very noteworthy role in two papers by
A.~Haar (\emph{Mathematische Zeitschrift}, vol.~31, no.~5, pp.~769--798, 1930, and
another appearing there in 1931).

\originalnote{2}The Hermitian (bounded; cf. A, Introduction, 1) operators $A,B$ are
called commuting if $AB=BA$.

\originalnote{3}For the description of Hilbert space, cf. E, Chapter~I (and also
Introduction I--IV and Appendix~I).  For operators, cf. E, Chapter~II and Appendix~II,
\S1.  Everything needed is also collected once more in A, Introduction, 1.  The ring
of all bounded operators mentioned immediately below is introduced in A,
Introduction, 2.  We denote points of Hilbert space by
$f,g,\ldots,\varphi,\psi,\ldots$; bounded operators by $A,B,\ldots$; and complex
numbers by $a,\ldots$ (as in E and A).

\originalnote{4}In matrix terminology, $A^*$ is the conjugate transpose of $A$; cf.
E, Appendix~II, \S1, or A, end of Introduction, 1.

\originalnote{5}The definition is in A, Chapter~II, \S1, Definition~1.  Cf. also A,
Chapter~I, \S3, the supplement to Theorem~5.

\originalnote{6}Cf. A, Chapter~I, \S1, Definition~2.

\originalnote{7}Cf. A, end of Chapter~III, \S2.

\originalnote{8}That is, $A=A^*$, respectively $A$ and $A^*$ commute.  The latter is
the criterion for $\mathbf R(A)$ to be Abelian, that is, to consist entirely of
mutually commuting operators.  Cf. A, Chapter~III, beginning of \S1 and \S2.

\originalnote{9}Cf. A, toward the end of Chapter~II, \S1.

\originalnote{10}Chapter~III, \S2, Theorem~10.

\originalnote{11}Cf. W.~H.~Young, \emph{Proceedings of the Cambridge Philosophical
Society}, vol.~14, pp.~520--529, 1908; H.~Lebesgue, thesis, \emph{Annali di
Matematica}, ser.~3, vol.~7, 1902.

\originalnote{12}$g(u)$ need no longer be measurable!

\originalnote{13}Since $\psi(b)$ is right-continuous, it certainly assumes
at the left endpoint of the interval the value assumed in the interior; at the right
endpoint, however, it does so only if it is continuous there.  That is, this happens
when the value $\psi(b)=a$ is not the beginning of an interval of constancy of
$\varphi(a)$.  Thus the indicated interval is to be understood as including its left
endpoint in any case, but as including its right endpoint only in the case just
described.

% Source PDF page 34 / journal p. 224
\originalnote{14}Here it is essential that $\mathfrak{P}_b$ included all the $I_n$,
and hence that $\mathfrak{P}'_b$ did also; had $\mathfrak{P}_b$ been disjoint from
some of the $I_n$, $\mathfrak{P}'_b$ would have had to do neither of these things.

\originalnote{15}Since one may put $\Phi(a)=a$, ordinary measurability follows from
this $\Phi$-measurability.

\originalnote{16}For a single $f_1(x)$, namely $f_1(x)=x$ (but with the second Baire
class in place of the third), this is a familiar theorem of Vitali.  The essential
point here is precisely the generalization to arbitrary measurable functions
$f_1(x),f_2(x),\ldots$.  For Vitali's theorem, cf. \emph{Rendiconti Lombardi},
vol.~38, p.~599 (1905).

\originalnote{17}For example, let $H(x)$ be of the second class, let
$H_1(x),H_2(x),\ldots$ be of the first class, and suppose $H_n(x)\to H(x)$.  If
$H(x)$ is to be made constant in the intervals $I_1,I_2,\ldots$, we accomplish this
by making $H_1(x)$ constant only in $I_1$, $H_2(x)$ constant only in $I_1,I_2$, and
so on.  In our case ($H(x)=G(x)$ and, after the alteration, $H(x)=\overline{G}(x)$),
the $I_n$ include their left endpoints at most; whether they include their right
endpoints as well is uncertain (cf. Note~\noteref{13}).

\originalnote{18}By dividing the number line into finitely many pieces, each
containing exactly one such point, one sees that it is enough to deal with a single
such point.  A trivial change of variable further permits the interval of definition
to be normalized to $-1<x<1$ and the point in question to $x=0$.  Let the sequence of
functions in question now be $f_1(x),f_2(x),\ldots$.  Define
\[
 \overline{f}_n(x)=f_n(x)
 \quad\text{for } |x|\geq\frac1n \text{ and for }x=0,
\]
and define $\overline{f}_n(x)$ linearly on each of
\[
 -\frac1n\leq x\leq0,
 \qquad
 0\leq x\leq\frac1n.
\]
Then $\overline{f}_1(x),\overline{f}_2(x),\ldots$ are continuous and have the same
limits as $f_1(x),f_2(x),\ldots$.

\originalnote{19}$f_{m,n}(x)\to g_m(x)$ as $n\to\infty$, and
$g_m(x)\to\overline{G}(x)$ as $m\to\infty$.

\originalnote{20}If $f(x)$ is right-continuous, then all the functions
\[
 \overline{f}_n(y)=n\int_y^{,y+1/n}f(x)\,dx
\]
are continuous, and $\overline{f}_n(x)\to f(x)$.

\originalnote{21}This condition is fulfilled if $\mathfrak{M}$ contains all numbers
in the interval of definition of $\varphi(x)$ except those in its intervals of
constancy, for it then contains all the values of $\varphi(y)$.  Incidentally, it is
easy to see that the $y$ for which $\psi(y)$ does not belong to $\mathfrak{M}$ are the
following.  For every $x$ lying outside $\mathfrak{M}$ (but in the interval of
definition of $\varphi(x)$), take the number $y=\varphi(x)$ when $x$ is a point of
continuity of $\varphi(x)$, and take the whole ``jump interval''
\[
 \lim_{\substack{x'\to x\\x'<x}}\varphi(x')
 \leq y\leq
 \varphi(x)
 =\lim_{\substack{x'\to x\\x'>x}}\varphi(x')
\]
when $x$ is a point of discontinuity.  Thus the complement of $\mathfrak{M}$ within
the interval of definition must satisfy the following requirements: it must contain
no point of discontinuity, and its image under $y=\varphi(x)$ must be a Lebesgue null
set.

\originalnote{22}$\mathfrak{M}$ is to be $\Phi$-measurable; that is, the function
equal to $1$ on $\mathfrak{M}$ and to $0$ on its complement is to be
$\Phi$-measurable.  Equivalently, for every monotone nondecreasing function
$\varphi(x)$ right-continuous, the set of $x$ for which
$\varphi(x)\in\mathfrak{M}$ is to be Lebesgue measurable.

\originalnote{23}As is well known, one may form the upper and lower bounds of the real
parts and of the imaginary parts of all expressions
\[
 \varphi(x_1)-\varphi(x'_1)+\cdots+
 \varphi(x_n)-\varphi(x'_n),
 \qquad
 x_1<x'_1<\cdots<x_n<x'_n<x,
\]
where $x$ is fixed and everything else is variable, and set them, respectively, equal
to
\[
 \psi_1(x)-\Re\varphi(-\infty),
 \qquad -\psi_2(x),
 \qquad \psi_3(x)-\Im\varphi(-\infty),
 \qquad -\psi_4(x).
\]
These are then the desired functions.

\originalnote{24}Hilbert, \emph{G\"ottinger Nachrichten}, 1906, especially
pp.~189--190, Theorem~II.  Our notation differs somewhat from Hilbert's; it is the
notation used in E and A.  Cf. in particular E, pp.~91--92, Definition~17 and
Theorem~36, and p.~114, Appendix~II, Theorem~9*.

% Source PDF page 35 / journal p. 225
\originalnote{25}The integrand, $\lambda$, is continuous, and
$(E(\lambda)f,g)$ is of bounded variation.  The latter follows from
\[
 \Re(E(\lambda)f,g)
 =\left(E(\lambda)\frac{f+g}{2},\frac{f+g}{2}\right)
  -\left(E(\lambda)\frac{f-g}{2},\frac{f-g}{2}\right),
\]
with an analogous expression for $\Im(E(\lambda)f,g)$ obtained by using $ig$ in
place of $g$.  Thus $(E(\lambda)f,g)$ is a linear combination of four terms
$(E(\lambda)h,h)$ with coefficients $\pm1,\pm i$; and
\[
 (E(\lambda)h,h)=| E(\lambda)h|^2
\]
is monotone nondecreasing in $\lambda$ and bounded (cf. E, p.~114).

\originalnote{26}Cf. A, p.~418, Appendix~I.

\originalnote{27}Cf. E, p.~91, Definition~17.  $E\leq F$ means
$| Ef|\leq| Ff|$ for every $f$; $E_n\to E$ means
$E_nf\to Ef$ for every $f$ (in the sense of ``strong'' convergence).

\originalnote{28}Since $\mathfrak{S}$ is always bounded, this is certainly the case
if $F(x)$ is bounded in every finite interval.  For example, $F(x)=x$ has this
property.

\originalnote{29}For $E\leq F$ one has in general
\[
 | Ef|^2-| Ff|^2
 =-|(F-E)f|^2;
\]
cf. E, p.~78.

\originalnote{30}$E(\lambda_n)-E(\lambda_{n-1})=E$ is a projection operator, and
for such operators one has in general
\[
 |(Ef,g)|\leq| Ef|\,| Eg|
\]
(cf. E, p.~93).

\originalnote{31}Cf. the citation and proof in E, p.~94, Note~52.

\originalnote{32}For real-valued $F(x)$ it follows in particular that $F(A)$ is
Hermitian.  In view of what was said in 1, it even suffices that $F(x)$ be real on the
spectrum $\mathfrak{S}$ of $A$.

\originalnote{33}Since $B=H(A)$ in $G(B)$ must be Hermitian, and $y=H(x)$ in $G(y)$
must be real, $H(x)$ must always be real (cf. Note~\noteref{32}).  By
\hyperref[thm:2]{Theorem~2}, $F(x)$ will be $\Phi$-measurable if $G(x)$ and $H(x)$
are.

\originalnote{34}Exactly as in Note~\noteref{21}, one sees that this is equivalent to
the following.  After the complement of $\mathfrak{M}$ (whose general element is
denoted by $\xi$) has been subjected to the mapping $x=\varrho^{-1}(\xi)$, it may
contain no point of discontinuity of $\Phi_f(x)$; hence the complement itself may
contain no point of discontinuity of $| E(\xi)f|^2$ (which is true for every
$f$ if it is disjoint from the point spectrum $\mathfrak{P}$ of $A$).  Furthermore,
its image under $x=\varrho^{-1}(\xi)$ and $y=\Phi_f(x)$---that is, its image under
\[
 y=| E(\xi)f|^2
\]
---must be a Lebesgue null set.  Since this function is always constant in the
complementary intervals of the spectrum, only those parts of the complement of
$\mathfrak{M}$ that lie in the spectrum $\mathfrak{S}$ of $A$ need be taken into
account (cf. also 1).  Thus, if $\mathfrak{M}$ contains all of $\mathfrak{S}$, every
condition is fulfilled.

\originalnote{35}Cf. A, p.~383.

\originalnote{36}Cf. E, p.~65, condition C.

\originalnote{37}Put $F_n(A)f=f_n$ and $F(A)f=\overline{f}$.  It must then be shown
that weak convergence $f_n\to\overline{f}$ together with
$| f_n|\to|\overline{f}|$ implies strong convergence
$f_n\to\overline{f}$ (for the terminology, cf., e.g., A, pp.~378--379).  Indeed,
\begin{align*}
 | f_n-\overline{f}|^2
 &= (f_n-\overline{f},f_n-\overline{f}) \\
 &= | f_n|^2+|\overline{f}|^2
    -2\Re(f_n,\overline{f}) \\
 &\longrightarrow
 |\overline{f}|^2+|\overline{f}|^2
    -2\Re(\overline{f},\overline{f})=0,
\end{align*}
so $f_n\to\overline{f}$ strongly, as asserted.

\originalnote{38}The premise of 2, h) is supplied by the premise of $\gamma)$.  It is
true that in 2, h) (for given $A$, $F_n(x)$, and $F(x)$) the inference is made from
all $f$ to all $f$; but it is immediately clear that one may just as well infer from
one given $f$ (without regard to the others) to that same $f$.

\originalnote{39}It reads as follows.  Let $\eta>0$, and let
$\mathfrak{M}_1,\mathfrak{M}_2,\ldots,\mathfrak{M}$ be Lebesgue-measurable sets such
that the measures of $\mathfrak{M}_1,\mathfrak{M}_2,\ldots$ are all at least $\eta$,
the measure of $\mathfrak{M}$ is finite, and
$\mathfrak{M}_1,\mathfrak{M}_2,\ldots$ are all subsets of $\mathfrak{M}$.  Then the
set of all points belonging to infinitely many of
$\mathfrak{M}_1,\mathfrak{M}_2,\ldots$ has measure at least $\eta$.  Cf., e.g., de la
Vall\'ee Poussin, \emph{Cours d'analyse infinit\'esimale}, ed.~1/2
(Louvain--Paris, 1909), pp.~68--69.

\originalnote{40}Assuming the truth of the so-called continuum hypothesis (according
to which every uncountable set of numbers has the cardinality of the continuum), one
can give an example

% Source PDF page 36 / journal p. 226
of a $\Phi$-measurable function possessing an inverse that is not
$\Phi$-measurable.  We shall not discuss it further here.  (If
$\Phi$-measurability is replaced by ordinary measurability, a counterexample
certainly exists.)

\originalnote{41}It is easy to prove that one could equally well restrict attention
to functions $f_1(u),f_2(u),\ldots$ that are continuous everywhere.

\originalnote{42}We now define the zeroth class to be the totality of all functions
that actually depend on only finitely many variables and are continuous in those
variables.  We interpret the limiting process in the broader sense just described.
For example, the function defined by
\[
 F(x_1,x_2,\ldots)=
 \begin{cases}
   x_1+x_2+\cdots, & \text{if this series converges},\\
   0, & \text{otherwise}
 \end{cases}
\]
will belong to the first class.

\originalnote{43}This is a kind of Peano curve, used for such purposes by Steinhaus
(\emph{Studia Mathematica}, vol.~II, pp.~24--27, 1930).  It is based on the dyadic
expansion
\[
 z=0.\xi_1\xi_2\cdots
   =\sum_{\nu=1}^{\infty}\frac{\xi_\nu}{2^\nu},
 \qquad \xi_\nu\in\{0,1\},
\]
with infinitely many zeros, which is uniquely possible for every $0\leq z<1$, and is
defined by
\[
 \begin{gathered}
 y_1=0.\eta^{(1)}_1\eta^{(1)}_2\cdots,
 \quad\ldots,\quad
 y_n=0.\eta^{(n)}_1\eta^{(n)}_2\cdots,\\
 y=0.\eta^{(1)}_1\cdots\eta^{(n)}_1
       \eta^{(1)}_2\cdots\eta^{(n)}_2\cdots.
 \end{gathered}
\]
It maps the whole cube
$0\leq y_1<1,\ldots,0\leq y_n<1$ one-to-one onto all points
$y=0.\xi_1\xi_2\cdots$ of the interval $0\leq y<1$ for which each of the sequences
\[
 \xi_k,\xi_{k+n},\xi_{k+2n},\ldots
 \qquad(k=1,\ldots,n)
\]
contains infinitely many zeros---that is, onto all points except for a Lebesgue null
set.  That this mapping has the stated measure-preserving property is readily proved
by an argument analogous to Steinhaus's at the place cited.

\originalnote{44}This mapping too is most simply described by means of dyadic
expansions.  If the number of variables $x_1,x_2,\ldots$ is finite, say $m$, put
\[
 \begin{gathered}
 x_1=0.\xi^{(1)}_1\xi^{(1)}_2\cdots,
 \quad\ldots,\quad
 x_m=0.\xi^{(m)}_1\xi^{(m)}_2\cdots,\\
 x=0.\xi^{(1)}_1\cdots\xi^{(m)}_1
      \xi^{(1)}_2\cdots\xi^{(m)}_2\cdots;
 \end{gathered}
\]
if their number is infinite, put
\[
 \begin{gathered}
 x_1=0.\xi^{(1)}_1\xi^{(1)}_2\cdots,
 \qquad
 x_2=0.\xi^{(2)}_1\xi^{(2)}_2\cdots,
 \qquad\ldots,\\
 x=0.\xi^{(1)}_1\xi^{(2)}_1\xi^{(1)}_2
      \xi^{(3)}_1\xi^{(2)}_2\xi^{(1)}_3\cdots.
 \end{gathered}
\]
All the properties stated in the text are readily verified.

\subsection*{Translator's Notes}

\translatornote{A}On source p.~217, the two occurrences following the initial
definitions of $\alpha$ and the $\beta_i$ interchange those symbols.  The translation
uses the mathematically consistent convention stated in the initial definitions:
substitute the $\beta_i(y)$, and then replace $y$ by $\alpha(y_1,\ldots,y_n)$.

\translatornote{B}In the displayed definition on source p.~219, the condition in the
$H'_D(x)$ branch is printed as $|H(x)|\leq D$.  Since this branch defines $H'_D$ and
the surrounding argument uses $H'$, the translation restores the evidently intended
prime: $|H'(x)|\leq D$.

\translatornote{C}In the definitions of $\mathfrak M_v(\epsilon)$ on source p.~209
and $\mathfrak M_{m,n}(\epsilon)$ on source p.~211, the displayed inequalities are
printed as $\geq\epsilon$.  Every subsequent use of these sets in the proofs requires
the complementary inequality $<\epsilon$ (nonmembership is then equivalent to
$\geq\epsilon$).  The translation uses $<\epsilon$ in both definitions.

\end{document}
